% Chapter 10. All quoted numbers come from chapters/ch10-numbers.tex, written by code/ch10/run_all.py.
% Generated by code/ch10/run_all.py. Do not edit by hand.
\newcommand{\cvAprCarry}{0.03}
\newcommand{\cvAprSpread}{75}
\newcommand{\cvBackSpread}{\ensuremath{-}17.3}
\newcommand{\cvBackSpreadAbs}{17.3}
\newcommand{\cvBackWhen}{April 1996}
\newcommand{\cvLastWhen}{April 2024}
\newcommand{\cvMedSpread}{0.1}
\newcommand{\cvMedWhen}{March 1999}
\newcommand{\ihBBill}{0.76}
\newcommand{\ihBCre}{0.86}
\newcommand{\ihBFarm}{1.16}
\newcommand{\ihBHpi}{0.76}
\newcommand{\ihBMkt}{\ensuremath{-}0.75}
\newcommand{\ihBOil}{11.57}
\newcommand{\ihBReit}{0.41}
\newcommand{\ihBTimber}{0.95}
\newcommand{\ihCorrBill}{\ensuremath{-}0.25}
\newcommand{\ihCorrCre}{\ensuremath{-}0.05}
\newcommand{\ihCorrFarm}{0.04}
\newcommand{\ihCorrHpi}{\ensuremath{-}0.15}
\newcommand{\ihCorrMkt}{\ensuremath{-}0.27}
\newcommand{\ihCorrOil}{0.48}
\newcommand{\ihCorrReit}{\ensuremath{-}0.05}
\newcommand{\ihCorrTimber}{\ensuremath{-}0.01}
\newcommand{\ihEBill}{1.08}
\newcommand{\ihECre}{0.79}
\newcommand{\ihEFarm}{0.62}
\newcommand{\ihEHpi}{0.68}
\newcommand{\ihEMkt}{\ensuremath{-}0.39}
\newcommand{\ihEOil}{9.77}
\newcommand{\ihEReit}{\ensuremath{-}1.31}
\newcommand{\ihETimber}{0.75}
\newcommand{\ihEseBill}{0.03}
\newcommand{\ihEseCre}{0.36}
\newcommand{\ihEseFarm}{0.45}
\newcommand{\ihEseHpi}{0.28}
\newcommand{\ihEseMkt}{1.05}
\newcommand{\ihEseOil}{3.30}
\newcommand{\ihEseReit}{2.06}
\newcommand{\ihEseTimber}{0.60}
\newcommand{\ihFirst}{1971}
\newcommand{\ihHBillFive}{1.03}
\newcommand{\ihHBillTen}{1.27}
\newcommand{\ihHBillThree}{0.91}
\newcommand{\ihHCreFive}{0.72}
\newcommand{\ihHCreTen}{0.71}
\newcommand{\ihHCreThree}{0.82}
\newcommand{\ihHFarmFive}{0.68}
\newcommand{\ihHFarmTen}{0.40}
\newcommand{\ihHFarmThree}{0.92}
\newcommand{\ihHHpiFive}{0.76}
\newcommand{\ihHHpiTen}{0.86}
\newcommand{\ihHHpiThree}{0.80}
\newcommand{\ihHMktFive}{0.10}
\newcommand{\ihHMktTen}{0.42}
\newcommand{\ihHMktThree}{\ensuremath{-}0.22}
\newcommand{\ihHReitFive}{\ensuremath{-}1.76}
\newcommand{\ihHReitTen}{\ensuremath{-}2.30}
\newcommand{\ihHReitThree}{\ensuremath{-}1.40}
\newcommand{\ihHTimberFive}{2.10}
\newcommand{\ihHTimberTen}{2.83}
\newcommand{\ihHTimberThree}{1.49}
\newcommand{\ihHnFive}{51}
\newcommand{\ihHnTen}{46}
\newcommand{\ihHnThree}{53}
\newcommand{\ihHsBillFive}{0.20}
\newcommand{\ihHsBillTen}{0.25}
\newcommand{\ihHsBillThree}{0.16}
\newcommand{\ihHsCreFive}{0.30}
\newcommand{\ihHsCreTen}{0.23}
\newcommand{\ihHsCreThree}{0.32}
\newcommand{\ihHsFarmFive}{0.31}
\newcommand{\ihHsFarmTen}{0.25}
\newcommand{\ihHsFarmThree}{0.38}
\newcommand{\ihHsHpiFive}{0.25}
\newcommand{\ihHsHpiTen}{0.26}
\newcommand{\ihHsHpiThree}{0.28}
\newcommand{\ihHsMktFive}{0.55}
\newcommand{\ihHsMktTen}{0.49}
\newcommand{\ihHsMktThree}{0.79}
\newcommand{\ihHsReitFive}{0.92}
\newcommand{\ihHsReitTen}{1.10}
\newcommand{\ihHsReitThree}{1.42}
\newcommand{\ihHsTimberFive}{1.03}
\newcommand{\ihHsTimberTen}{1.15}
\newcommand{\ihHsTimberThree}{0.67}
\newcommand{\ihInflMean}{3.8}
\newcommand{\ihInflSd}{2.8}
\newcommand{\ihInflTwo}{13.1}
\newcommand{\ihLast}{2025}
\newcommand{\ihNBill}{55}
\newcommand{\ihNCre}{55}
\newcommand{\ihNFarm}{55}
\newcommand{\ihNHpi}{50}
\newcommand{\ihNMkt}{55}
\newcommand{\ihNOil}{39}
\newcommand{\ihNReit}{29}
\newcommand{\ihNTimber}{44}
\newcommand{\ihRealBill}{0.6}
\newcommand{\ihRealCre}{0.3}
\newcommand{\ihRealFarm}{\ensuremath{-}0.8}
\newcommand{\ihRealHpi}{1.3}
\newcommand{\ihRealMkt}{6.8}
\newcommand{\ihRealOil}{0.5}
\newcommand{\ihRealReit}{3.0}
\newcommand{\ihRealTimber}{\ensuremath{-}0.7}
\newcommand{\ihRealVolBill}{2.6}
\newcommand{\ihSeBill}{0.14}
\newcommand{\ihSeCre}{0.28}
\newcommand{\ihSeFarm}{0.51}
\newcommand{\ihSeHpi}{0.29}
\newcommand{\ihSeMkt}{1.25}
\newcommand{\ihSeOil}{3.17}
\newcommand{\ihSeReit}{2.55}
\newcommand{\ihSeTimber}{0.51}
\newcommand{\ihStartBill}{1971}
\newcommand{\ihStartCre}{1971}
\newcommand{\ihStartFarm}{1971}
\newcommand{\ihStartHpi}{1976}
\newcommand{\ihStartMkt}{1971}
\newcommand{\ihStartOil}{1987}
\newcommand{\ihStartReit}{1997}
\newcommand{\ihStartTimber}{1982}
\newcommand{\ihTwoBill}{\ensuremath{-}11.1}
\newcommand{\ihTwoBillAbs}{11.1}
\newcommand{\ihTwoCre}{\ensuremath{-}1.8}
\newcommand{\ihTwoCreAbs}{1.8}
\newcommand{\ihTwoFarm}{31.8}
\newcommand{\ihTwoFarmAbs}{31.8}
\newcommand{\ihTwoHpi}{13.7}
\newcommand{\ihTwoHpiAbs}{13.7}
\newcommand{\ihTwoMkt}{\ensuremath{-}13.9}
\newcommand{\ihTwoMktAbs}{13.9}
\newcommand{\ihTwoOil}{35.5}
\newcommand{\ihTwoOilAbs}{35.5}
\newcommand{\ihTwoReit}{\ensuremath{-}16.8}
\newcommand{\ihTwoReitAbs}{16.8}
\newcommand{\ihTwoTimber}{\ensuremath{-}1.2}
\newcommand{\ihTwoTimberAbs}{1.2}
\newcommand{\ihUBill}{0.17}
\newcommand{\ihUCre}{0.99}
\newcommand{\ihUFarm}{2.14}
\newcommand{\ihUHpi}{1.00}
\newcommand{\ihUMkt}{\ensuremath{-}1.40}
\newcommand{\ihUOil}{13.00}
\newcommand{\ihUReit}{0.33}
\newcommand{\ihUTimber}{1.22}
\newcommand{\ihUseBill}{0.04}
\newcommand{\ihUseCre}{0.30}
\newcommand{\ihUseFarm}{0.46}
\newcommand{\ihUseHpi}{0.29}
\newcommand{\ihUseMkt}{1.52}
\newcommand{\ihUseOil}{4.57}
\newcommand{\ihUseReit}{2.56}
\newcommand{\ihUseTimber}{0.52}
\newcommand{\ihVarCut}{10}
\newcommand{\ihVolBill}{2.6}
\newcommand{\ihVolCre}{8.2}
\newcommand{\ihVolFarm}{10.1}
\newcommand{\ihVolHpi}{3.9}
\newcommand{\ihVolMkt}{17.7}
\newcommand{\ihVolOil}{32.2}
\newcommand{\ihVolReit}{17.4}
\newcommand{\ihVolTimber}{6.8}
\newcommand{\ihWcre}{0.09}
\newcommand{\lcApprPeak}{79}
\newcommand{\lcApprShareEighty}{0}
\newcommand{\lcApprWorst}{1989Q4}
\newcommand{\lcFirst}{1996Q2}
\newcommand{\lcGapMax}{11.0}
\newcommand{\lcLastStart}{2021Q2}
\newcommand{\lcLateEnd}{50}
\newcommand{\lcLateFirst}{2016Q1}
\newcommand{\lcLtvZero}{65}
\newcommand{\lcMedPeak}{65}
\newcommand{\lcN}{101}
\newcommand{\lcShareEighty}{13}
\newcommand{\lcShareHundred}{2}
\newcommand{\lcShareRefi}{14}
\newcommand{\lcSpread}{1.80}
\newcommand{\lcWorst}{2007Q3}
\newcommand{\lcWorstEnd}{80}
\newcommand{\lcWorstEndLtv}{80}
\newcommand{\lcWorstEndVintage}{2007Q3}
\newcommand{\lcWorstPeak}{102}
\newcommand{\lcWorstPx}{\ensuremath{-}26}
\newcommand{\lcWorstPxAbs}{26}
\newcommand{\lcWorstRate}{6.80}
\newcommand{\lrLtvNine}{44}
\newcommand{\lrLtvThree}{65}
\newcommand{\lrSwitchHigh}{6.36}
\newcommand{\lrSwitchLow}{3.73}
\newcommand{\lsMeanFive}{0.76}
\newcommand{\lsMeanOne}{\ensuremath{-}0.02}
\newcommand{\lsMeanTen}{0.84}
\newcommand{\lsN}{10,000}
\newcommand{\lsRejFive}{10}
\newcommand{\lsRejOne}{8}
\newcommand{\lsRejTen}{19}
\newcommand{\lsSdFive}{0.08}
\newcommand{\lsSdOne}{0.13}
\newcommand{\lsSdTen}{0.08}
\newcommand{\lsSeFive}{0.07}
\newcommand{\lsSeOne}{0.12}
\newcommand{\lsSeTen}{0.05}
\newcommand{\lsT}{55}
\newcommand{\lvAmort}{25}
\newcommand{\lvBal}{76.9}
\newcommand{\lvCapSd}{1.00}
\newcommand{\lvCashYield}{1.27}
\newcommand{\lvCfYield}{5.00}
\newcommand{\lvDebtYield}{9.9}
\newcommand{\lvDscrAtMax}{1.08}
\newcommand{\lvDscrMin}{1.25}
\newcommand{\lvGsd}{2.5}
\newcommand{\lvIlMean}{7.54}
\newcommand{\lvIlQ}{2.51}
\newcommand{\lvIlSd}{3.12}
\newcommand{\lvIuMean}{7.10}
\newcommand{\lvIuQ}{4.58}
\newcommand{\lvIuSd}{1.66}
\newcommand{\lvK}{7.92}
\newcommand{\lvLevRet}{9.10}
\newcommand{\lvLevRetEff}{8.87}
\newcommand{\lvLoan}{14.8}
\newcommand{\lvLtv}{56.1}
\newcommand{\lvLtvMax}{65}
\newcommand{\lvNsim}{20,000}
\newcommand{\lvPBelow}{36.1}
\newcommand{\lvPLoss}{0.8}
\newcommand{\lvPLossU}{0.0}
\newcommand{\lvRate}{6.25}
\newcommand{\lvRateEff}{6.43}
\newcommand{\lvRateSwitch}{4.74}
\newcommand{\lvSdRatio}{1.87}
\newcommand{\lvTerm}{10}
\newcommand{\lvVolMult}{2.28}
\newcommand{\mgGrowthValue}{385}
\newcommand{\mgRent}{80}
\newcommand{\pfCap}{5.56}
\newcommand{\pfCapNoCapex}{5.00}
\newcommand{\pfCapTenth}{6.11}
\newcommand{\pfCapUp}{\ensuremath{-}15.3}
\newcommand{\pfCapUpAbs}{15.3}
\newcommand{\pfCapex}{10}
\newcommand{\pfCapexAmt}{0.146}
\newcommand{\pfCf}{1.316}
\newcommand{\pfDcfGap}{\ensuremath{-}3.9}
\newcommand{\pfDcfGapAbs}{3.9}
\newcommand{\pfEgi}{2.282}
\newcommand{\pfElast}{\ensuremath{-}0.621}
\newcommand{\pfExit}{5.81}
\newcommand{\pfG}{2.5}
\newcommand{\pfGpr}{2.40}
\newcommand{\pfIrr}{6.98}
\newcommand{\pfIrrSame}{7.50}
\newcommand{\pfNoi}{1.462}
\newcommand{\pfOpex}{0.82}
\newcommand{\pfOpexShare}{35.9}
\newcommand{\pfOther}{0.05}
\newcommand{\pfR}{7.5}
\newcommand{\pfSaleCost}{2}
\newcommand{\pfSaleShare}{62.1}
\newcommand{\pfT}{10}
\newcommand{\pfTenthExact}{5.6}
\newcommand{\pfTenthLog}{5.9}
\newcommand{\pfTenthLogVal}{5.7}
\newcommand{\pfUpHundred}{\ensuremath{-}9.5}
\newcommand{\pfUpHundredAbs}{9.5}
\newcommand{\pfVac}{7}
\newcommand{\pfVacAmt}{0.168}
\newcommand{\pfValue}{26.3}
\newcommand{\pfValueDcf}{25.3}
\newcommand{\pfValueNoCapex}{29.2}
\newcommand{\rtBetaCre}{\ensuremath{-}0.03}
\newcommand{\rtBetaCreDim}{0.32}
\newcommand{\rtBetaReit}{0.72}
\newcommand{\rtBetaReitDim}{0.83}
\newcommand{\rtBuniReit}{\ensuremath{-}0.02}
\newcommand{\rtBzeroReit}{0.00}
\newcommand{\rtCorrCreMkt}{\ensuremath{-}0.06}
\newcommand{\rtCorrCreReit}{\ensuremath{-}0.05}
\newcommand{\rtCorrReitMkt}{0.59}
\newcommand{\rtCorrY}{0.34}
\newcommand{\rtCorrYLag}{0.63}
\newcommand{\rtFirst}{1997Q1}
\newcommand{\rtLagOne}{\ensuremath{-}0.02}
\newcommand{\rtLast}{2026Q2}
\newcommand{\rtLeadFour}{0.22}
\newcommand{\rtLeadOne}{0.41}
\newcommand{\rtLeadTwo}{0.32}
\newcommand{\rtN}{118}
\newcommand{\rtNY}{29}
\newcommand{\rtSumReit}{0.43}
\newcommand{\rtUnlevFive}{10.9}
\newcommand{\rtUnlevFour}{13.0}
\newcommand{\rtUnlevThree}{15.2}
\newcommand{\rtUvolCre}{9.4}
\newcommand{\rtVolCre}{7.6}
\newcommand{\rtVolCreY}{10.3}
\newcommand{\rtVolMkt}{17.7}
\newcommand{\rtVolReit}{21.7}
\newcommand{\rtVolReitY}{17.4}
\newcommand{\rtYFirst}{1997}
\newcommand{\rtYLast}{2025}
\newcommand{\scBigRes}{\ensuremath{-}39}
\newcommand{\scBigYear}{2009}
\newcommand{\scResSd}{8.2}
\newcommand{\scYears}{55}
\newcommand{\sgAprCy}{\ensuremath{-}223}
\newcommand{\sgAprCyAbs}{223}
\newcommand{\sgAprDev}{10.3}
\newcommand{\sgAprFour}{29.17}
\newcommand{\sgAprOne}{16.70}
\newcommand{\sgBackFirst}{January 1990}
\newcommand{\sgBackShare}{49}
\newcommand{\sgCorrGas}{\ensuremath{-}0.43}
\newcommand{\sgCorrOil}{\ensuremath{-}0.46}
\newcommand{\sgCurveLast}{April 2024}
\newcommand{\sgCyMean}{2.4}
\newcommand{\sgDevHigh}{25.4}
\newcommand{\sgDevLow}{\ensuremath{-}7.0}
\newcommand{\sgGapMean}{9.8}
\newcommand{\sgGapSe}{4.5}
\newcommand{\sgGapT}{2.2}
\newcommand{\sgGasCurveFirst}{January 1994}
\newcommand{\sgGasFirst}{January 1994}
\newcommand{\sgGasLast}{April 2024}
\newcommand{\sgGasMissing}{June 2014}
\newcommand{\sgJanFirst}{1997}
\newcommand{\sgJanLast}{2026}
\newcommand{\sgJanMean}{8.9}
\newcommand{\sgJanN}{29}
\newcommand{\sgJanPos}{66}
\newcommand{\sgJanSd}{40.1}
\newcommand{\sgLockFirst}{1997}
\newcommand{\sgLockLast}{2023}
\newcommand{\sgLockMean}{15.3}
\newcommand{\sgLockMin}{2.0}
\newcommand{\sgLockN}{27}
\newcommand{\sgLockPos}{100}
\newcommand{\sgLockSd}{11.6}
\newcommand{\sgMaxContango}{223}
\newcommand{\sgNGas}{363}
\newcommand{\sgNGasSpan}{364}
\newcommand{\sgNOil}{196}
\newcommand{\sgOilFirst}{January 2008}
\newcommand{\sgOilLast}{April 2024}
\newcommand{\sgQHigh}{\ensuremath{-}17.6}
\newcommand{\sgQLow}{11.0}
\newcommand{\sgQuadOil}{1.9}
\newcommand{\sgRealMean}{5.5}
\newcommand{\sgRealPos}{56}
\newcommand{\sgRealSd}{26.1}
\newcommand{\sgSlopeDec}{\ensuremath{-}9}
\newcommand{\sgSlopeGas}{\ensuremath{-}1.09}
\newcommand{\sgSlopeJan}{\ensuremath{-}29}
\newcommand{\sgSlopeOil}{\ensuremath{-}0.91}
\newcommand{\sgSlopeSep}{92}
\newcommand{\sgStockLast}{July 2026}
\newcommand{\smAcfFourHpi}{0.55}
\newcommand{\smAcfFourPost}{0.04}
\newcommand{\smAcfFourPre}{0.58}
\newcommand{\smAcfFourReit}{\ensuremath{-}0.10}
\newcommand{\smArHpi}{0.60}
\newcommand{\smArHpiX}{0.605}
\newcommand{\smArPost}{0.24}
\newcommand{\smArPostX}{0.236}
\newcommand{\smArPre}{0.85}
\newcommand{\smArPreX}{0.849}
\newcommand{\smArReit}{0.07}
\newcommand{\smArReitX}{0.066}
\newcommand{\smCorrHalf}{0.26}
\newcommand{\smCorrShrink}{0.53}
\newcommand{\smCreFall}{\ensuremath{-}39}
\newcommand{\smCreFallAbs}{39}
\newcommand{\smCrePeak}{2007Q3}
\newcommand{\smCreTrough}{2009Q4}
\newcommand{\smCumGap}{\ensuremath{-}11.9}
\newcommand{\smCumRep}{107.4}
\newcommand{\smCumTrue}{95.5}
\newcommand{\smHpiFirst}{1975Q2}
\newcommand{\smLast}{2026Q2}
\newcommand{\smMeanHpi}{4.8}
\newcommand{\smMeanPost}{4.1}
\newcommand{\smMeanPre}{4.3}
\newcommand{\smMeanReit}{5.9}
\newcommand{\smNHpi}{205}
\newcommand{\smNPost}{121}
\newcommand{\smNPre}{103}
\newcommand{\smNReit}{118}
\newcommand{\smPostFirst}{1996Q2}
\newcommand{\smPreFirst}{1970Q2}
\newcommand{\smPreLast}{1995Q4}
\newcommand{\smRatioHpi}{2.02}
\newcommand{\smRatioPost}{1.28}
\newcommand{\smRatioPre}{3.43}
\newcommand{\smReitBreak}{73}
\newcommand{\smReitFall}{\ensuremath{-}69}
\newcommand{\smReitFallAbs}{69}
\newcommand{\smReitFirst}{1997Q1}
\newcommand{\smReitPeak}{2007Q1}
\newcommand{\smReitTrough}{2009Q1}
\newcommand{\smTheoryHpi}{2.02}
\newcommand{\smTheoryPost}{1.27}
\newcommand{\smTheoryPre}{3.50}
\newcommand{\smUarHpi}{\ensuremath{-}0.06}
\newcommand{\smUarPost}{\ensuremath{-}0.01}
\newcommand{\smUarPre}{\ensuremath{-}0.05}
\newcommand{\smUvolHpi}{5.5}
\newcommand{\smUvolPost}{9.6}
\newcommand{\smUvolPre}{11.3}
\newcommand{\smVolHpi}{2.7}
\newcommand{\smVolPost}{7.5}
\newcommand{\smVolPre}{3.3}
\newcommand{\smVolReit}{21.7}
\newcommand{\ssGapEnd}{\ensuremath{-}4.6}
\newcommand{\ssGapMax}{10.4}
\newcommand{\ssQ}{40}
\newcommand{\ssVolFactor}{0.29}
\newcommand{\ssVolModel}{2.4}
\newcommand{\ssVolRep}{1.8}
\newcommand{\ssVolTrue}{8.3}
\newcommand{\tbA}{400}
\newcommand{\tbC}{3}
\newcommand{\tbCfOne}{100}
\newcommand{\tbConFifty}{1,963}
\newcommand{\tbConNinetyNine}{2,191}
\newcommand{\tbConThirty}{1,610}
\newcommand{\tbCost}{1,500}
\newcommand{\tbCpiAnn}{3.82}
\newcommand{\tbDurFifty}{17.2}
\newcommand{\tbDurNinetyNine}{22.4}
\newcommand{\tbDurPerp}{23.8}
\newcommand{\tbDurThirty}{12.4}
\newcommand{\tbFFour}{25.3}
\newcommand{\tbFPriceBoth}{27.1}
\newcommand{\tbFPriceGrowth}{26.0}
\newcommand{\tbFSix}{22.6}
\newcommand{\tbFThree}{27.1}
\newcommand{\tbFirstFarm}{1970}
\newcommand{\tbFirstTimber}{1981}
\newcommand{\tbGcon}{2.5}
\newcommand{\tbGrowthFour}{5.05}
\newcommand{\tbGrowthSix}{6.23}
\newcommand{\tbGrowthThree}{4.42}
\newcommand{\tbK}{0.06}
\newcommand{\tbLastYear}{2025}
\newcommand{\tbLevFour}{1,994}
\newcommand{\tbLevSFour}{1,933}
\newcommand{\tbLevSSix}{250}
\newcommand{\tbLevSThree}{3,711}
\newcommand{\tbLevSix}{251}
\newcommand{\tbLevThree}{3,914}
\newcommand{\tbLumberPeak}{164}
\newcommand{\tbLumberPeakDate}{May 2021}
\newcommand{\tbMai}{31.7}
\newcommand{\tbPerp}{2,222}
\newcommand{\tbPrice}{40}
\newcommand{\tbRcon}{7}
\newcommand{\tbRcross}{3.2}
\newcommand{\tbRealFarm}{\ensuremath{-}0.82}
\newcommand{\tbRealFarmAbs}{0.82}
\newcommand{\tbRealLumber}{\ensuremath{-}0.23}
\newcommand{\tbRealLumberAbs}{0.23}
\newcommand{\tbRealTimber}{\ensuremath{-}0.73}
\newcommand{\tbRealTimberAbs}{0.73}
\newcommand{\tbRealVolFarm}{10.1}
\newcommand{\tbRealVolLumber}{14.2}
\newcommand{\tbRealVolTimber}{6.8}
\newcommand{\tbSFour}{28.4}
\newcommand{\tbSSix}{23.1}
\newcommand{\tbSThree}{32.4}
\newcommand{\tbShareFifty}{88.3}
\newcommand{\tbShareNinetyNine}{98.6}
\newcommand{\tbShareThirty}{72.4}
\newcommand{\tbShockPerp}{\ensuremath{-}18.2}
\newcommand{\tbShockPerpAbs}{18.2}
\newcommand{\tbShockThirty}{\ensuremath{-}10.6}
\newcommand{\tbShockThirtyAbs}{10.6}
\newcommand{\tbVolMai}{247}
\newcommand{\tbYearsTimber}{44}
\newcommand{\xaCap}{6.52}
\newcommand{\xaNoi}{740}
\newcommand{\xaValue}{11.35}
\newcommand{\xbApprox}{\ensuremath{-}5.9}
\newcommand{\xbApproxPct}{\ensuremath{-}5.7}
\newcommand{\xbDrop}{\ensuremath{-}5.6}
\newcommand{\xbShare}{62.0}
\newcommand{\xcIrr}{8}
\newcommand{\xcIrrHigh}{7.03}
\newcommand{\xdDscr}{1.26}
\newcommand{\xdK}{6.813}
\newcommand{\xeK}{7.732}
\newcommand{\xeLoan}{19.90}
\newcommand{\xeLtv}{62.2}
\newcommand{\xeLtvLoan}{22.4}
\newcommand{\xfBal}{76.4}
\newcommand{\xgEqVol}{25}
\newcommand{\xgRet}{9.25}
\newcommand{\xgRetLow}{1.75}
\newcommand{\xhAcf}{0.24}
\newcommand{\xhVol}{4.2}
\newcommand{\xiOne}{2.25}
\newcommand{\xiThree}{\ensuremath{-}3.25}
\newcommand{\xiTwo}{\ensuremath{-}1.00}
\newcommand{\xjSum}{0.9375}
\newcommand{\xkF}{23.9}
\newcommand{\xkLev}{918}
\newcommand{\xkS}{25.4}
\newcommand{\xlDur}{10.6}
\newcommand{\xlShare}{69.4}
\newcommand{\xlV}{694}
\newcommand{\xmBound}{72.13}
\newcommand{\xmCy}{1.16}
\newcommand{\xmGross}{3.16}
\newcommand{\xnCov}{\ensuremath{-3.6\times 10^{-4}}}
\newcommand{\xnW}{0.054}
\newcommand{\xoLast}{0.038}
\newcommand{\xoSeasOne}{0.748}
\newcommand{\xoSeasThird}{0.408}
\newcommand{\xoV}{0.463}

\begin{refsection}
\setcounter{chapter}{9}
\chapter{Real Assets and Commercial Real Estate}
\label{ch:realassets}

An office building, a warehouse, a forest, a farm, a toll road and a tank of crude oil have one thing
in common: each is a physical asset whose value comes from what it produces or what it saves its owner,
and the physical assets do not trade continuously like listed shares, although exchange-traded futures
provide exposure to commodities. Their prices are found in private sales,
appraisals and futures markets, their returns are measured by indexes that are themselves estimates,
and many, commercial property and infrastructure above all, are bought with substantial borrowing. This chapter is about how such assets are valued, how
their returns are measured, and what they did for an investor worried about inflation.

The chapter starts with a building: how rents become net operating income, how net operating income
becomes a value through a cap rate or a discounted cash flow, and how sensitive that value is to the
price at which the building is sold. Leverage comes next, with the two tests a lender applies, the
loan they allow and what the loan does to the owner's return and risk; a case study marks loans
written in each quarter from \lcFirst{} to \lcLastStart{} to the Federal Reserve's commercial property index. The third section is about
measurement. Appraisals smooth the true path of values, and a short model shows by how much the
smoothing hides volatility and correlation and how to undo it; the Federal Reserve's commercial
property index, which switched from appraisals to transaction prices in 1996, tests the model. Listed
real estate investment trusts (REITs) are then set against private property. The last three sections
turn to timberland, farmland and infrastructure, where time and biological growth enter the value, to
stored commodities, where the futures curve reveals the value of holding inventory and lets the owner of
a tank fix what it will earn, and to inflation: which of these assets kept their real value when prices
rose, over a year and over a decade, and how much of a hedge each one gave.

\begin{keybox}[How to read this chapter]
The numbers in the text are computed by nine scripts in \texttt{code/ch10/} (\texttt{proforma},
\texttt{leverage}, \texttt{cycle}, \texttt{smoothing}, \texttt{reits}, \texttt{timber},
\texttt{storage}, \texttt{inflation} and \texttt{exercises}), and the tests in
\texttt{code/ch10/test\_*.py} check the library against simulation, finite differences and numerical
optimisation. An independent R script and a GNU Octave script recompute the main results by other
routes, and \texttt{code/ch10/README.md} lists what each of them checks.

The property, REIT and house price series are quarterly from the Financial Accounts of the United
States and the Federal Housing Finance Agency, and the price indexes, bill and bond yields and spot
energy prices are monthly, all published on FRED \parencite{fredch10}; they run to the second quarter
or to September 2026. The futures curves and inventories are from the US Energy Information
Administration \parencite{eiach10}. US equity returns are from the Kenneth R. French Data Library
\parencite{frenchlib}. The building, the loan terms, the forest and the concession are illustrations,
and are labelled as such where they appear.

Random draws come from NumPy's default generator with base seed \texttt{3817} and a stream number fixed
in each script. The seed is an arbitrary label that makes the runs repeatable; it carries no date or
other meaning.

Notation. $\mathit{NOI}_t$ is net operating income in year $t$, $V$ a value, $c$ a cap rate, $r$ a
required return or discount rate, $g$ a growth rate and $\kappa$ the share of NOI spent on capital items.
For a loan, $\lambda$ is the loan-to-value ratio, $k$ the mortgage constant and $i$ the interest rate.
In \cref{sec10:smoothing}, $r_t$ is a true return, $r^*_t$ the reported one and $\alpha$ the smoothing
weight. $F(t,T)$ is a futures price for delivery at $T$ and $S_t$ the spot price; $u$ is a storage cost
and $y$ a convenience yield, both continuous rates. $\pi_t$ is inflation over a year. Returns in the smoothing model
and the regressions are log returns; the valuation, loan and leverage formulas use simple annual rates,
and mortgage payments compound monthly. In the forest and storage sections $r$ is continuously
compounded. Volatilities are annual. Six short listings in the text, kept in
\texttt{code/ch10/snippets/}, run from the repository root and print the numbers they are placed beside.
\end{keybox}

\FloatBarrier
% =====================================================================================
\section{From rents to value}
\label{sec10:value}

A commercial building earns rent from its tenants and pays the costs of running it. \emph{Net
operating income} (NOI) is what is left: potential rent at full occupancy, less vacancy and rent that
is not collected, plus other income such as parking, less the operating expenses that the leases do
not pass on to the tenants. NOI is measured before debt service, income tax and depreciation, so it
describes the property and not the way it is financed or owned. It is also measured before capital
spending. Roofs, lifts and heating plants wear out, and each new lease costs a commission to the broker
and an allowance for the tenant's fit out. Those outlays are lumpy, and a valuation spreads them as a
reserve, a share $\kappa$ of NOI, so that the cash the property returns to its owner is $(1-\kappa)\mathit{NOI}$.

\Cref{tab10:proforma} sets out a pro forma for a stylised suburban office building. The numbers are
illustrations chosen to be of a realistic order; they are not those of any building.

\begin{table}[t]
  \centering
  \small
  \caption{Pro forma for the first year of a stylised office building, in millions of dollars. Value is
  next year's cash flow capitalised at $r-g$, with $r=\pfR\%$ and $g=\pfG\%$.}
  \label{tab10:proforma}
  \begin{tabular}{@{}lr@{}}
    \toprule
    Gross potential rent & \pfGpr\\
    less vacancy and credit loss (\pfVac\%) & \pfVacAmt\\
    plus parking and other income & \pfOther\\
    Effective gross income & \pfEgi\\
    less operating expenses not recovered & \pfOpex\\
    Net operating income & \pfNoi\\
    less capital reserve (\pfCapex\% of NOI) & \pfCapexAmt\\
    Cash flow before debt service & \pfCf\\
    \addlinespace
    Value at a cap rate of \pfCap\% & \pfValue\\
    \bottomrule
  \end{tabular}
\end{table}

\subsection{The cap rate}

The \emph{capitalisation rate} or cap rate of a property is next year's NOI divided by its value or
price. It is reported for most sales, and it plays the role that the earnings yield plays for a share. Its
link to returns is the growing perpetuity of \textcite{gordon1956}, adjusted for the capital reserve.

\begin{proposition}[Cap rate, required return and growth]\label{prop10:cap}
Suppose NOI grows at a constant rate $g$, the owner spends a constant share $\kappa<1$ of NOI on capital
items, and cash flows are discounted at $r>g$. Then the value of the property is
\begin{equation}\label{eq10:gordon}
  V = \frac{(1-\kappa)\,\mathit{NOI}_1}{r-g},
  \qquad\text{so}\qquad
  c = \frac{\mathit{NOI}_1}{V} = \frac{r-g}{1-\kappa}.
\end{equation}
\end{proposition}

\begin{proof}
The cash flow in year $t$ is $(1-\kappa)\mathit{NOI}_1(1+g)^{t-1}$, and its present value summed over
$t\ge1$ is a geometric series with ratio $(1+g)/(1+r)<1$:
$\sum_{t\ge1}(1-\kappa)\mathit{NOI}_1(1+g)^{t-1}(1+r)^{-t}=(1-\kappa)\mathit{NOI}_1/(r-g)$. Dividing
$\mathit{NOI}_1$ by $V$ gives the cap rate.
\end{proof}

With a required return of \pfR\%, growth of \pfG\% and a reserve of \pfCapex\%, the cap rate is
\pfCap\% and the building is worth \pfValue{} million dollars. Leaving the reserve out would give a cap rate of
\pfCapNoCapex\% and a value of \pfValueNoCapex{} million, an error of more than a tenth. The comparison
says something about quoted cap rates. Two buildings that sell at the same cap rate need not offer the
same return: one may need a new roof, and the other may have rents below the market that will rise
when its leases end. A cap rate is a price, and \cref{eq10:gordon} turns it into a return only once
the growth and the capital needs are known.

\subsection{Discounted cash flow and the exit cap rate}

Investors who buy a building usually plan to sell it, and value it by a discounted cash flow over a
holding period of $T$ years, with a sale at the end at a price set by an \emph{exit} cap rate $c_T$
applied to the NOI of the year after:
\begin{equation}\label{eq10:dcf}
  V_{\mathrm{DCF}} = \sum_{t=1}^{T}\frac{(1-\kappa)\mathit{NOI}_t}{(1+r)^t}
  + \frac{(1-s)\,\mathit{NOI}_{T+1}/c_T}{(1+r)^T},
\end{equation}
where $s$ is the cost of selling as a share of the price. The two methods should agree when their
assumptions agree, and the next result says when they do.

\begin{proposition}[Consistency of the two methods]\label{prop10:dcf}
Let NOI grow at $g$, and let $c=(r-g)/(1-\kappa)$ as in \cref{prop10:cap}. If $c_T=c$ and $s=0$, then
$V_{\mathrm{DCF}}=V$ for every $T$. A buyer who pays $V$ and later sells at the same cap rate earns an
internal rate of return of exactly $r=(1-\kappa)c+g$, whatever the holding period.
\end{proposition}

\begin{proof}
At $c_T=c$ the sale price is $\mathit{NOI}_{T+1}/c=(1-\kappa)\mathit{NOI}_{T+1}/(r-g)$, which by
\cref{prop10:cap} applied at year $T$ is the value at $T$ of the cash flows after $T$. The sum in
\cref{eq10:dcf} is the value of the first $T$ cash flows, so the total is the value of all of them, $V$.
For the second statement, the price $V$ equals the present value at rate $r$ of the cash flows and the
sale, so $r$ solves the IRR equation; the IRR is unique because the flows change sign once.
\end{proof}

The return on a property bought and sold at the same cap rate is its cash yield plus its growth, the
real estate form of the growth identity that \textcite{gordon1956} wrote for shares. Exit cap rates are, in practice, often
set above going-in rates, on the view that the building will be older when it is sold. With an exit
cap of \pfExit\%, a quarter of a point above the going-in rate, and selling costs of \pfSaleCost\%,
the DCF value is \pfValueDcf{} million, \pfDcfGapAbs\% below the value in \cref{tab10:proforma}. A buyer who
paid the full \pfValue{} million and then sold on those terms would earn \pfIrr\% and not \pfIrrSame\%.

The listing below computes these figures from the inputs of \cref{tab10:proforma}. Its DCF with the
exit cap set equal to the going-in cap returns the Gordon value of \pfValue{} million, a check that the
cash flows and the sale are dated consistently; the higher exit cap and the selling costs then give
\pfValueDcf{} million, and the internal rates of return at the Gordon price are \pfIrr\% and \pfIrrSame\%.
\Needspace*{28\baselineskip}
\lstinputlisting[language=Python]{code/ch10/snippets/dcf.py}

\begin{proposition}[Sensitivity to the exit cap rate]\label{prop10:exit}
The elasticity of $V_{\mathrm{DCF}}$ with respect to the exit cap rate is minus the share of the value
that comes from the sale,
\[
  \frac{\partial\log V_{\mathrm{DCF}}}{\partial\log c_T}
  = -\frac{(1-s)\mathit{NOI}_{T+1}/c_T}{(1+r)^T\,V_{\mathrm{DCF}}},
\]
and at $c_T=c$ and $s=0$ that share is $\bigl((1+g)/(1+r)\bigr)^{T}$.
\end{proposition}

\begin{proof}
Only the sale term of \cref{eq10:dcf} depends on $c_T$, and it is proportional to $1/c_T$, so its
derivative with respect to $\log c_T$ is minus itself. Dividing by $V_{\mathrm{DCF}}$ gives the first
formula. At $c_T=c$, the sale is worth the value at $T$ of all later cash flows,
$V(1+g)^T$, and its present value is $V(1+g)^T(1+r)^{-T}$, so the share is
$\bigl((1+g)/(1+r)\bigr)^T$.
\end{proof}

At \pfR\% and \pfG\% over \pfT{} years the sale carries \pfSaleShare\% of the value, so a rise of one
tenth in the exit cap rate, from \pfCap\% to \pfCapTenth\%, lowers the value by about \pfTenthLog\% in logs to first
order, $\pfSaleShare\%\times\ln1.1$, or \pfTenthLogVal\% in value, and by \pfTenthExact\% exactly. A full point higher exit cap lowers the value by \pfUpHundredAbs\%.
The same point added to the going-in cap rate would lower it by \pfCapUpAbs\%, because it applies to every year and not only to the sale. Panel (b) of
\cref{fig10:exitcap} shows the share falling with the holding period: at short horizons the DCF is
mostly a forecast of the exit cap rate, which is a forecast of the market's mood years ahead.

\begin{figure}[!tb]
  \centering
  \includegraphics[width=\textwidth]{ch10/fig01_exitcap}
  \caption{(a) DCF value of the building of \cref{tab10:proforma} against the exit cap rate, for holding
  periods of 5, 10 and 20 years; the dotted line is the going-in cap rate, where the three agree. (b)
  Share of the value that comes from the sale, $((1+g)/(1+r))^T$, at $r=\pfR\%$ and three growth rates.}
  \label{fig10:exitcap}
\end{figure}

\FloatBarrier
% =====================================================================================
\section{Leverage}
\label{sec10:leverage}

Most commercial property is bought partly with a mortgage. A commercial mortgage is usually written for
a \emph{term} of five to ten years with payments set as if the loan were repaid over a longer
\emph{amortisation} period, so that a balance is still owed at the end of the term and must be repaid
or refinanced. Two numbers describe the payments.

\begin{proposition}[Mortgage constant and balance]\label{prop10:loan}
A loan of one dollar at a nominal annual rate $i>0$ compounded monthly, repaid in equal monthly payments
over $n$ years, has the
annual debt service, or \emph{mortgage constant},
\begin{equation}\label{eq10:constant}
  k(i,n) = \frac{12\,j}{1-(1+j)^{-12n}}, \qquad j=i/12,
\end{equation}
and after $t$ years, with $12t$ a whole number of payments, the balance is
\begin{equation}\label{eq10:balance}
  B_t = \frac{(1+j)^{12n}-(1+j)^{12t}}{(1+j)^{12n}-1}
  = \frac{k(i,n)}{12}\cdot\frac{1-(1+j)^{-12(n-t)}}{j}.
\end{equation}
\end{proposition}

\begin{proof}
The monthly payment $p$ makes the present value of $12n$ payments at rate $j$ equal to one:
$p\,(1-(1+j)^{-12n})/j=1$, and $k=12p$. The balance evolves as $B_{m+1}=(1+j)B_m-p$ from $B_0=1$, whose
solution is $B_m=(1+j)^m-p\,((1+j)^m-1)/j$; substituting $p$ gives the first form. The second form is
the present value of the $12(n-t)$ payments still due, which equals the balance because the loan is
repaid exactly at $n$.
\end{proof}

A lender sizes a loan by two tests, which Chapter~6 met in a credit screen. The
\emph{loan-to-value ratio} (LTV) is the loan divided by the appraised value, and the \emph{debt service
coverage ratio} (DSCR) is NOI divided by the annual debt service. Chapter~6 showed how these ratios and
the debt yield tie together; the question here is the owner's, how large a loan the two tests allow.

\begin{proposition}[The two tests]\label{prop10:dscr}
For a property with value $V$ and cap rate $c=\mathit{NOI}_1/V$, a loan of $\lambda V$ with mortgage
constant $k$ has
\begin{equation}\label{eq10:dscr}
  \mathrm{DSCR} = \frac{c}{\lambda k}.
\end{equation}
With a highest LTV of $\bar\lambda$ and a lowest DSCR of $\underline D$, the largest loan is
$V\min\bigl(\bar\lambda,\;c/(\underline D\,k)\bigr)$. The DSCR test binds when $c<\underline D\,\bar\lambda\,k$,
that is, when the cap rate is low against the mortgage constant.
\end{proposition}

\begin{proof}
Debt service is $k\lambda V$ and NOI is $cV$, so their ratio is $c/(\lambda k)$, a form of the identities of
Chapter~6. The DSCR test
$c/(\lambda k)\ge\underline D$ holds if and only if $\lambda\le c/(\underline Dk)$, and the largest loan
meets the tighter of the two upper bounds on $\lambda$.
\end{proof}

At \lvRate\% over \lvAmort{} years the mortgage constant is \lvK\%, above the building's cap rate of
\pfCap\%. A loan at \lvLtvMax\% of value would have a DSCR of only \lvDscrAtMax, so a lender who asks for
\lvDscrMin{} lends \lvLoan{} million dollars, \lvLtv\% of value. The coverage test, not the value test,
sets the loan, and it does so for any rate above \lvRateSwitch\%. When rates rise faster than cap rates,
buyers can borrow less against the same building, which is one way higher rates reach property values.
The lender's own yardstick, the \emph{debt yield} of NOI over the loan, is \lvDebtYield\%.

\Cref{fig10:loanrate} draws the largest loan against the mortgage rate for this building and for
buildings priced at cap rates of 5\% and 6.5\%. While the value test binds, the loan is \lvLtvMax\% of
value whatever the rate; once the coverage test binds, the loan falls as the rate rises, to \lrLtvNine\%
of value for this building at 9\%. The kink moves with the cap rate, from \lrSwitchLow\% at a cap of 5\%
to \lrSwitchHigh\% at 6.5\%, so buildings bought at low cap rates are the first to lose borrowing
capacity when rates rise.

\begin{figure}[!tb]
  \centering
  \includegraphics[width=0.8\textwidth]{ch10/fig14_loanrate}
  \caption{Largest loan as a share of value, against the mortgage rate, for a lender who asks for a DSCR
  of at least \lvDscrMin{} and an LTV of at most \lvLtvMax\% on a loan amortised over \lvAmort{} years. The
  solid line is the building of \cref{tab10:proforma}, at a cap rate of \pfCap\%; the dotted line marks
  \lvRateSwitch\%, above which the coverage test sets its loan.}
  \label{fig10:loanrate}
\end{figure}

\begin{proposition}[Leverage, return and risk]\label{prop10:lever}
Over a single period, let the property return $R_A$ and let a loan of a share $\lambda<1$ of value,
with its principal fixed at the start and repaid with certainty at the end, cost $i$, fixed in advance. The return on the owner's equity is
\begin{equation}\label{eq10:lever}
  R_E = R_A + \frac{\lambda}{1-\lambda}\,(R_A-i),
  \qquad \sigma(R_E)=\frac{\sigma(R_A)}{1-\lambda}.
\end{equation}
Leverage raises the expected return on equity if and only if $\E R_A>i$.
\end{proposition}

\begin{proof}
The owner puts up $(1-\lambda)$ per unit of value and receives $R_A-\lambda i$ on it, so
$R_E=(R_A-\lambda i)/(1-\lambda)$, which rearranges to \cref{eq10:lever}. As $i$ is fixed,
$\sigma(R_E)=\sigma(R_A)/(1-\lambda)$, and $\E R_E-\E R_A=\lambda(\E R_A-i)/(1-\lambda)$ has the sign of
$\E R_A-i$.
\end{proof}

Whether leverage is positive depends on which return is set against the cost of debt. For the building
the expected return is \pfR\% against a rate of \lvRate\%, so leverage at \lvLtv\% raises the expected
equity return to \lvLevRet\% and multiplies its volatility by \lvVolMult. That is a one-period figure that sets the \pfR\% return against
the quoted \lvRate\%; paid monthly, the loan costs \lvRateEff\% a year, which gives \lvLevRetEff\%. The first year's cash yield
tells a different story: the cash flow over value is \lvCfYield\%, but after debt service the equity
earns only \lvCashYield\% in cash, because the mortgage constant, which includes repayment of principal,
is above the cap rate. The higher expected return comes from earning the property's \pfR\% on money
borrowed for less, and almost all of it arrives through growth in value, which is uncertain, and at the
sale, when the principal repaid along the way comes back as a smaller balance to pay off.

\Cref{fig10:levered} puts numbers on that uncertainty. On each of \lvNsim{} simulated paths the
building's NOI grows each year at \pfG\% plus a normal shock with a standard deviation of \lvGsd\%, and
the building is sold after \lvTerm{} years at an exit cap rate drawn around \pfExit\% with a standard
deviation of \lvCapSd{} point; the loan above is repaid from the sale at its balance of \lvBal\% of the
original amount, and each year's twelve payments are carried to the end of the year at the loan rate,
so the loan costs \lvRateEff\% a year. Across the paths the unlevered IRR averaged \lvIuMean\% with a standard deviation of \lvIuSd\%, and the
levered IRR \lvIlMean\% with \lvIlSd\%, \lvSdRatio{} times as much. The 5\% quantile fell from \lvIuQ\%
to \lvIlQ\%. Levered equity did worse than the property itself on \lvPBelow\% of the paths, exactly those
on which the property's IRR fell below the loan's \lvRateEff\%, and it lost money on \lvPLoss\% of them;
no unlevered path lost money. The levered mean is below the \lvLevRet\% of \cref{prop10:lever} because the
higher exit cap pulls the unlevered mean down to \lvIuMean\%, the loan costs \lvRateEff\% and not
\lvRate\%, and amortisation and rising value lower the leverage over the ten years.

\begin{figure}[!tb]
  \centering
  \includegraphics[width=0.86\textwidth]{ch10/fig02_levered}
  \caption{Ten year IRRs of the building of \cref{tab10:proforma}, bought at \pfValue{} million dollars,
  on \lvNsim{} simulated paths of rent growth and the exit cap rate, without debt and with a loan of
  \lvLtv\% of value at \lvRate\%. The dotted line is the loan's cost of \lvRateEff\% a year,
  the quoted \lvRate\% paid monthly.}
  \label{fig10:levered}
\end{figure}

\subsection{Case study: loan-to-value through the cycle}
\label{sec10:cycle}

Chapter~6 followed Canadian five year loans to their refinancing. The question here is a different one:
how far can a loan written at a prudent LTV drift while it is outstanding, when the value of the building
moves with the market? Take a stylised loan written at the start of each quarter from \lcFirst{} to
\lcLastStart, \lcN{} vintages in all, at \lcLtvZero\% of value, at the US 10 year Treasury yield of the
quarter's first month plus \lcSpread{} points, amortised over 25 years with a five year term. Mark the
building each quarter with the Federal Reserve's commercial property price index, the repeat-sales index
of \cref{sec10:smoothing}, so that the LTV after $t$ years is
\[
  \lambda_t = \lambda_0\,\frac{B_t}{P_t/P_0},
\]
with $B_t$ the balance of \cref{prop10:loan} and $P_t$ the index. The spread, the starting LTV and the
use of a national index for a single building are assumptions; a single building can move far more
than an index of many.

The calculation needs only the index, the Treasury yield and the balance formula, and the listing
below runs it from the data files with the chapter's own loaders. It prints the number of vintages, the
worst vintage with its highest LTV and the share of vintages whose LTV passed 80\% at some point, the
numbers discussed next.
\Needspace*{21\baselineskip}
\lstinputlisting[language=Python]{code/ch10/snippets/ltv_cycle.py}

\Cref{fig10:cycle} shows the result. For the median vintage the highest LTV over the term was \lcMedPeak\%,
no higher than at the start, because amortisation and rising prices pulled it down. But \lcShareEighty\% of
the vintages passed 80\% at some point and \lcShareHundred\% passed 100\%. The worst was the loan of
\lcWorst, written at \lcWorstRate\% at the peak of the index: its LTV reached \lcWorstPeak\% and was still
\lcWorstEnd\% at maturity, when the index stood \lcWorstPxAbs\% below its level at origination. Refinancing it at
\lcLtvZero\% would have needed new equity of \lcGapMax\% of the building's original value. In all,
\lcShareRefi\% of the vintages ended their term above the starting LTV, and the loans written from
\lcLateFirst{} on ended at about \lcLateEnd\% on average.

\begin{figure}[!tb]
  \centering
  \includegraphics[width=\textwidth]{ch10/fig03_cycle}
  \caption{(a) Loan-to-value ratio over the five year term of loans written at \lcLtvZero\% in four
  quarters, marked to the Z.1 commercial property price index. (b) Highest loan-to-value over the term
  and loan-to-value at maturity for loans written in each quarter from \lcFirst{} to \lcLastStart.}
  \label{fig10:cycle}
\end{figure}

The same calculation on the appraisal-era index tells a gentler story about the early 1990s, when US
commercial property went through a deep fall. Of the loans written from 1985 to 1990, none passed 80\%,
and the highest LTV, for the vintage of \lcApprWorst, was \lcApprPeak\%. Some of the difference is real, but
\cref{sec10:smoothing} shows how much an appraisal index smooths, and a lender who marked its book to
such an index would have seen a smoother path of collateral values. This comparison alone cannot say
how much of the gap came from smoothing and how much from different markets.

\FloatBarrier
% =====================================================================================
\section{Appraisal smoothing}
\label{sec10:smoothing}

The return on a share is measured from two prices at which the share traded. A building trades every
ten years or so, and the return on a portfolio of buildings in a quarter is usually measured from
appraisals: a valuer's estimate of what each building would sell for. Valuers work from the sales of
similar buildings in recent months and from the previous valuation, and both anchor the new estimate
to the past. An index built from appraisals therefore moves later and less than the values it
measures. Chapter~1 introduced the model of \textcite{geltner1993} for this,
\begin{equation}\label{eq10:smooth}
  r^*_t = \alpha\,r^*_{t-1} + (1-\alpha)\,r_t, \qquad 0\le\alpha<1,
\end{equation}
where $r_t$ is the true return in period $t$ and $r^*_t$ the reported one, and promised to return to it
here. Unrolling the recursion, $r^*_t=(1-\alpha)\sum_{j\ge0}\alpha^jr_{t-j}$: each true return enters the
reports spread over many later periods, with weights that add to one. Three earlier results carry over.
When the true returns are independent with variance $\sigma^2$, the reports keep the mean, have
variance $\sigma^2(1-\alpha)/(1+\alpha)$ and autocorrelation $\alpha^k$ at lag $k$ (Chapter~1); the slope
of the reports on the current market return is $(1-\alpha)$ times the true beta, and the sum of the slopes
on current and lagged market returns recovers it (Chapter~7); and, under the same model, smoothing
leaves the variance per period of long sums unchanged in the limit, though over a finite horizon the two
variances differ, and \cref{eq10:smooth} can be inverted to unsmooth a series (Chapter~8). Two more
consequences matter for property: what smoothing does to the return over a holding period, and what it
does to a correlation.

\begin{proposition}[Holding period returns and correlations under smoothing]\label{prop10:smooth}
Let $r^*_t$ follow \cref{eq10:smooth}.
\begin{enumerate}[label=(\alph*)]
  \item For any path and any start,
        $\sum_{t=1}^Kr^*_t=\sum_{t=1}^Kr_t-\frac{\alpha}{1-\alpha}\bigl(r^*_K-r^*_0\bigr)$.
  \item Let $r_t=a+\beta m_t+e_t$, where the market returns $m_t$ are independent with variance
        $\sigma_m^2$ and the $e_t$ are independent of each other and of the market, and let $r^*$ be
        stationary. Then the correlation of $r^*_t$ with $m_t$ is $\sqrt{1-\alpha^2}$ times the
        correlation of $r_t$ with $m_t$, and the correlation of $r^*_t$ with $m_{t-j}$ is $\alpha^j$
        times that of $r^*_t$ with $m_t$.
\end{enumerate}
\end{proposition}

\begin{proof}
(a) From \cref{eq10:smooth}, $(1-\alpha)r_t=r^*_t-\alpha r^*_{t-1}$. Summing over $t=1,\dots,K$ gives
$(1-\alpha)\sum r_t=(1-\alpha)\sum r^*_t+\alpha(r^*_K-r^*_0)$, which rearranges to the identity.
(b) In $r^*_t=(1-\alpha)\sum_i\alpha^i(a+\beta m_{t-i}+e_{t-i})$ only the term $i=j$ is correlated with
$m_{t-j}$, so $\Cov(r^*_t,m_{t-j})=(1-\alpha)\alpha^j\beta\sigma_m^2$, which gives the second statement.
With $\Var r^*_t=\sigma_r^2(1-\alpha)/(1+\alpha)$, the correlation at $j=0$ is
$(1-\alpha)\beta\sigma_m/\bigl(\sigma_r\sqrt{(1-\alpha)/(1+\alpha)}\bigr)=\sqrt{1-\alpha^2}\,\beta\sigma_m/\sigma_r$,
and $\beta\sigma_m/\sigma_r$ is the correlation of $r_t$ with $m_t$.
\end{proof}

Part (a) says that the difference between the cumulative true and reported log returns over a holding
period is an end-point term and not an accumulation of errors period by period: it is the change between
the first and last reports, scaled by $\alpha/(1-\alpha)$. Its size varies with those two reports, but it
does not build up as the period lengthens. What
the index gets wrong is the timing, and with it every statistic that depends on timing. Part (b) gives
the size of the error in a correlation, the input that matters most when property is mixed with shares
and bonds. At the $\alpha=\smArPreX$ of the appraisal-era index below, the factor is
\smCorrShrink, so a true correlation of 0.5 with the stock market would be reported as \smCorrHalf, and the asset would look like a much better diversifier than it is. The
lagged correlations decline geometrically, which is the pattern a lead and lag regression looks for.
\textcite{getmansky2004} apply the same reasoning to hedge funds with illiquid holdings.

\Cref{fig10:smoothsim} shows part (a) on simulated returns. True quarterly log returns are
drawn as independent normals with the mean and volatility of the repeat-sales index after 1996, and
smoothed at $\alpha=\smArPreX$ with the first report at the mean. Over \ssQ{} quarters the reported path
trails the true one without drifting away from it: the gap, true less reported, is never more than
\ssGapMax{} points in absolute value, ends at \ssGapEnd{} points, and equals the end-point term of part (a)
in every quarter. The reported returns have a volatility of \ssVolRep\% a year against \ssVolTrue\% for the
true ones. The model's factor $\sqrt{(1-\alpha)/(1+\alpha)}=\ssVolFactor$ would give \ssVolModel\%, and a
sample of \ssQ{} quarters of so persistent a series understates its variance, here by more than usual.

\begin{figure}[!tb]
  \centering
  \includegraphics[width=0.7\textwidth]{ch10/fig15_smoothsim}
  \caption{A simulation: cumulative true and reported log returns over \ssQ{} quarters, with true
  returns drawn as independent normals with the mean and volatility of the repeat-sales index and
  smoothed at $\alpha=\smArPreX$; random stream \texttt{[3817, 2]}.}
  \label{fig10:smoothsim}
\end{figure}

\paragraph{Unsmoothing}
If \cref{eq10:smooth} holds, the true return can be recovered from two reports, as in Chapters~1 and 8:
\begin{equation}\label{eq10:unsmooth}
  r_t = \frac{r^*_t-\alpha\,r^*_{t-1}}{1-\alpha},
\end{equation}
with $\alpha$ replaced by the estimated first autocorrelation. This is the correction of
\textcite{geltner1993}, who also showed that it does not need the market to be informationally efficient.
It divides the change in the reported series by $1-\alpha$, so it also amplifies any noise in the reports,
and it is only as good as the model. A valuer may smooth more in some years than in others, and an
index that mixes properties appraised in different quarters is smoothed by its construction as well.

\paragraph{The evidence}
The Federal Reserve's quarterly price index for commercial real estate in the Financial Accounts offers a
natural test. From 1996 it chains a transaction-based repeat-sales index, the CoStar U.S. Composite Index
excluding multifamily, value weighted; before 1996 it is built from appraisal-based price series for
offices, retail and industrial property. \Cref{tab10:smooth} sets the two periods side by side with the
FHFA all-transactions house price index, which mixes sales with appraisals made for refinancing, and with
the price return of equity REITs from the same accounts. The REIT series starts in \smReitFirst{} because
its revaluation in the last quarter of 1996, equal to \smReitBreak\% of the previous quarter's market
value, is a break in the series and not a return. The two property periods differ in market conditions
and coverage as well as in method, so their contrast is suggestive and not a controlled estimate of
smoothing.

\begin{table}[t]
  \centering
  \small
  \caption{Quarterly log price returns, annualised: mean, volatility (twice the quarterly standard
  deviation), first and fourth order autocorrelation, and the volatility after unsmoothing with
  \cref{eq10:unsmooth} at the estimated $\alpha$. Through \smLast. REITs are not unsmoothed, since their
  $\rho_1$ is not significant.}
  \label{tab10:smooth}
  \setlength{\tabcolsep}{4pt}
  \begin{tabular}{@{}lrrrrrr@{}}
    \toprule
    & Quarters & Mean & Volatility & $\rho_1$ & $\rho_4$ & Unsmoothed\\
    \midrule
    Property, appraisals, \smPreFirst{} to \smPreLast & \smNPre & \smMeanPre & \smVolPre & \smArPre & \smAcfFourPre & \smUvolPre\\
    Property, repeat sales, from \smPostFirst & \smNPost & \smMeanPost & \smVolPost & \smArPost & \smAcfFourPost & \smUvolPost\\
    Houses, FHFA, from \smHpiFirst & \smNHpi & \smMeanHpi & \smVolHpi & \smArHpi & \smAcfFourHpi & \smUvolHpi\\
    Equity REITs, price, from \smReitFirst & \smNReit & \smMeanReit & \smVolReit & \smArReit & \smAcfFourReit & --\\
    \bottomrule
  \end{tabular}
\end{table}

The appraisal-era statistics are consistent with what the model predicts. The first autocorrelation was
\smArPre, the fourth still \smAcfFourPre, and the volatility only \smVolPre\% a year. After the switch to
transaction prices the first autocorrelation fell to \smArPost{} and the fourth to \smAcfFourPost, and
the volatility more than doubled to \smVolPost\%. The mean return hardly changed, \smMeanPre\% a year
before and \smMeanPost\% after. Unsmoothing the
appraisal-era index at $\alpha=\smArPreX$ multiplies its volatility by \smRatioPre, close to the
$\sqrt{(1+\alpha)/(1-\alpha)}=\smTheoryPre$ that the model gives, to \smUvolPre\%, and leaves residual
autocorrelation of \smUarPre; unsmoothed volatilities depend on the smoothing model. The transaction index still has some persistence, perhaps because sales
close months after their prices are agreed, and unsmoothing it lifts its volatility to \smUvolPost\%. The
house price index, with its appraisals, sits in between. REIT prices, set continuously on an exchange, show no
reliable autocorrelation at all. \Cref{fig10:acf} draws the four autocorrelation functions.

The R listing below estimates $\alpha$ for the appraisal era, unsmooths the index by
\cref{eq10:unsmooth} and prints the reported and unsmoothed volatilities, the ratio of \smRatioPre{}
against \smTheoryPre, and the residual autocorrelation of \smUarPre.
\Needspace*{14\baselineskip}
\lstinputlisting[language=R]{code/ch10/snippets/unsmooth.R}

\begin{figure}[!tb]
  \centering
  \includegraphics[width=0.92\textwidth]{ch10/fig05_acf}
  \caption{Autocorrelations of quarterly price returns at lags of one to eight quarters, for the Z.1
  commercial property index before 1996 (appraisals) and after (repeat sales), the FHFA house price
  index and equity REITs. The shaded band is $\pm2/\sqrt{n}$ for the repeat-sales sample.}
  \label{fig10:acf}
\end{figure}

\Cref{fig10:levels} shows the timing in the financial crisis. The property index peaked in \smCrePeak{}
and fell \smCreFallAbs\% to its trough in \smCreTrough. Equity REIT prices peaked two quarters earlier, in
\smReitPeak, and fell \smReitFallAbs\% to a trough in \smReitTrough, three quarters before the property
index reached its own, by which time they had recovered much of the fall. The timing is consistent with faster price discovery in the
stock market, though the two indexes do not hold the same buildings.

\begin{figure}[!tb]
  \centering
  \includegraphics[width=0.92\textwidth]{ch10/fig04_levels}
  \caption{Price levels of commercial property (Z.1 index), equity REITs (cumulative revaluation over market
  value) and houses (FHFA all-transactions index), \smReitFirst{} = 100, log scale, through \smLast.}
  \label{fig10:levels}
\end{figure}

\FloatBarrier
% =====================================================================================
\section{REITs and private property}
\label{sec10:reits}

A real estate investment trust is a company that owns, and usually operates, income-producing property
and passes most of its income to its shareholders. In the United States the tax code sets the terms
\parencite{usc856}. A REIT must have at least 100 beneficial owners and must not be closely held; at the
end of each quarter at least 75\% of its assets must be real estate, cash and government securities; at
least 75\% of its gross income must come from rents, interest on mortgages and similar sources; and it
must pay out as dividends at least 90\% of its REIT taxable income, a tax measure with statutory
adjustments and not its cash flow. In return a qualifying REIT can deduct the dividends it pays, which
can remove federal corporate income tax on the income it distributes. Equity REITs own buildings, mortgage REITs own loans, and the shares of most
listed REITs trade on an exchange like any other share. An investor in an equity REIT therefore owns a
levered, managed and liquid claim on a portfolio of buildings whose own values are set in the private
market.

REIT returns and private property returns differ for three reasons that can be measured and one that
cannot easily be. REITs borrow, and their equity is a levered claim on the buildings. Private returns
are smoothed, by appraisals or by the lags of a transaction index. REITs and the private indexes
hold different mixes of property types. The fourth
reason is that a REIT is priced in the stock market, alongside every other share, and its price moves
with the market's appetite for risk as well as with the value of its buildings. The first is a short
calculation.

\begin{proposition}[Removing leverage]\label{prop10:unlever}
Over a period, let a REIT's assets be financed by equity and by debt that is a share $\lambda$ of the
market value of the assets at its start, with the three returns measured over the same period, so that the asset return is $R_A=(1-\lambda)R_E+\lambda R_D$. Then, for any return $m$,
\[
  \beta_A=(1-\lambda)\beta_E+\lambda\beta_D,\qquad
  \sigma_A^2=(1-\lambda)^2\sigma_E^2+2\lambda(1-\lambda)\Cov(R_E,R_D)+\lambda^2\sigma_D^2,
\]
where the betas are measured against $m$. If the debt return is known in advance, $\sigma_A=(1-\lambda)\sigma_E$
and the correlation of $R_A$ with $m$ equals that of $R_E$.
\end{proposition}

\begin{proof}
Covariance with $m$ is linear, which gives the beta, and the variance of a weighted sum gives the second
formula. With $R_D$ known, $\Cov(R_E,R_D)=\sigma_D^2=0$, so $\sigma_A=(1-\lambda)\sigma_E$, and
$\Cov(R_A,m)=(1-\lambda)\Cov(R_E,m)$, so the factor $(1-\lambda)$ cancels in the correlation.
\end{proof}

The proposition is \cref{prop10:lever} read backwards, and it has a consequence that is easy to miss:
leverage changes volatility and beta but, with safe debt, leaves the correlation with the stock market
where it was. That holds in the model; in practice refinancing, changing leverage, the repricing of debt and distress can also
change a REIT's sensitivity to the stock market.

\paragraph{The evidence}
Over the \rtN{} quarters from \rtFirst{} to \rtLast, the price return of equity REITs had a volatility of
\rtVolReit\% a year against \rtVolCre\% for the transaction-based property index (\smVolPost\% in
\cref{tab10:smooth}, whose sample starts three quarters earlier), and a correlation of
\rtCorrReitMkt{} with the US stock market against \rtCorrCreMkt{} for the index. The contemporaneous
correlation between REITs and the property index was \rtCorrCreReit, which would suggest two unrelated
assets. The lagged correlations say otherwise (\cref{fig10:leadlag}). The property index this quarter
had a correlation of \rtLeadOne{} with REIT returns one quarter earlier, \rtLeadTwo{} two quarters earlier
and \rtLeadFour{} four quarters earlier, and \rtLagOne{} with REIT returns one quarter later. REIT returns lead
the private index, and not the other way round. These correlations are descriptive: serial correlation,
changes in the index's composition and common shocks stop the largest lagged correlation being read as
a causal link or a stable forecast. A regression of the property index on current and four
lagged REIT returns gives slopes that sum to \rtSumReit, of which the current quarter contributes \rtBzeroReit; a regression on the current quarter alone gives a slope of \rtBuniReit.
The same lags appear against the stock market: the property index has a beta of \rtBetaCre{} on the current
market return and a Dimson beta of \rtBetaCreDim{} with four lags, while the REIT beta rises only from
\rtBetaReit{} to \rtBetaReitDim.

\begin{figure}[!tb]
  \centering
  \includegraphics[width=\textwidth]{ch10/fig06_leadlag}
  \caption{(a) Correlation of the quarterly return of the Z.1 commercial property index with the equity
  REIT price return $k$ quarters earlier, \rtFirst{} to \rtLast; positive $k$ means REITs lead. The shaded
  band is $\pm2/\sqrt{n}$. (b) Calendar year price returns of the property index against those of equity
  REITs, \rtYFirst{} to \rtYLast, with the least squares line.}
  \label{fig10:leadlag}
\end{figure}

At the annual horizon the two series come closer. Over the \rtNY{} calendar years from \rtYFirst{} to
\rtYLast{} the correlation of annual returns was \rtCorrY, and the correlation of the property index with
REIT returns in the previous year was \rtCorrYLag. The annual volatilities were \rtVolReitY\% for REITs and
\rtVolCreY\% for the index. Back in quarterly data, removing leverage with \cref{prop10:unlever}, whose
debt return is known in advance, at debt shares of 30\%, 40\% or 50\% of assets, a range chosen for
illustration, brings the REIT volatility of \rtVolReit\% down to \rtUnlevThree\%, \rtUnlevFour\% or
\rtUnlevFive\%, and unsmoothing the property index at its estimated $\alpha$ lifts its \rtVolCre\% to
\rtUvolCre\%. These show what the two corrections do; they are not estimates of the volatility of REIT
assets. The gap between \rtUvolCre\% and the unlevered REIT figures is the part that
neither correction removes. Part of it is the volatility that comes with a price set every day in the
stock market, where a REIT share moves with the market's mood, and part may come from differences in
property mix, operating leverage and valuation that these two corrections do not touch. \textcite{pagliari2005} restate public and private real estate indexes for property-type mix,
leverage and appraisal smoothing, and find that the restated means and volatilities are not statistically
different, a long-run agreement that the annual figures here also suggest.

% =====================================================================================
\section{Timberland, farmland and infrastructure}
\label{sec10:timber}

Three other real assets add something that a building does not have. A forest grows, so its owner chooses
when to harvest, and waiting is itself an investment. A farm earns a rent tied to the price of what it
grows. An infrastructure concession, a toll road or an airport, often comes with a fixed life, at the end
of which the asset returns to the state.

\subsection{The rotation of a forest}

Let a stand of trees planted on bare land have volume $Q(T)$ per hectare at age $T$, sold at a stumpage
price $p$ per cubic metre, net of harvesting costs, and replanted at a cost $C$. If the owner harvests and
replants every $T$ years forever, the value of the bare land at a continuously compounded rate $r$ is the
\emph{land expectation value} of \textcite{faustmann1995},
\begin{equation}\label{eq10:faustmann}
  \mathit{LEV}(T)=\frac{p\,Q(T)\,e^{-rT}-C}{1-e^{-rT}} .
\end{equation}

\begin{proposition}[The Faustmann rotation]\label{prop10:faustmann}
If $Q$ is differentiable and the rotation $T^*>0$ maximises \cref{eq10:faustmann}, then
\begin{equation}\label{eq10:foc}
  p\,Q'(T^*) = r\,p\,Q(T^*) + r\,\mathit{LEV}(T^*).
\end{equation}
If $\mathit{LEV}(T^*)>0$, and the ratio $Q'(T)/Q(T)$ is decreasing in $T$, then $T^*$ is shorter than the rotation that maximises the value of a single harvest,
$pQ(T)e^{-rT}$, which solves $Q'(T)=rQ(T)$.
\end{proposition}

\begin{proof}
Write $\mathit{LEV}(T)(1-e^{-rT})=pQ(T)e^{-rT}-C$ and differentiate at an interior maximum, where
$\mathit{LEV}'=0$: $\mathit{LEV}\cdot re^{-rT}=pQ'e^{-rT}-rpQe^{-rT}$. Dividing by $e^{-rT}$ gives
\cref{eq10:foc}. With $\mathit{LEV}>0$, \cref{eq10:foc} gives $Q'(T^*)/Q(T^*)>r$, and the single harvest
optimum $T_1$ has $Q'(T_1)/Q(T_1)=r$. As $Q'/Q$ is decreasing, $T^*<T_1$.
\end{proof}

\Cref{eq10:foc} is the optimal harvesting rule for a renewable asset under the model's assumptions. The owner waits one more year while the value
added by a year's growth, $pQ'$, exceeds the return the timber would earn if cut and sold, $rpQ$, plus
the rent on the land, $r\,\mathit{LEV}$, which the standing trees occupy. The land rent is what makes the
forest rotation shorter than the rotation for a single harvest, and in the example below it is also
shorter than the rotation that maximises average growth $Q(T)/T$, which ignores interest altogether.
The single harvest rotation is shorter than that only when $r$ exceeds $1/T$ at the average-growth
optimum, here \tbRcross\%, so at 3\% it is longer. \textcite{samuelson1976}
gives the history of the argument and of the errors made before Faustmann.

For an illustration take the growth curve of \textcite{richards1959} in the form
$Q(T)=\tbA\,(1-e^{-\tbK\,T})^{\tbC}$ cubic metres per hectare, a stumpage price of \tbPrice{} dollars per
cubic metre and a planting cost of \tbCost{} dollars per hectare. Average growth is greatest at
\tbMai{} years. At $r=3\%$ the Faustmann rotation is \tbFThree{} years and the single harvest rotation
\tbSThree, and the land is worth \tbLevThree{} dollars per hectare; at 4\% they are \tbFFour{} and \tbSFour{}
years, with land worth \tbLevFour{}; at 6\%, \tbFSix{} and \tbSSix{} years, and the land is worth only
\tbLevSix{} (\cref{fig10:rotation}). At the optimum the stand grows at \tbGrowthFour\% a year when $r=4\%$,
above the discount rate, as \cref{eq10:foc} requires. A timber price that rises at 1\% a year in real terms
works on the harvests like a fall of one point in the discount rate, while the planting cost is still
discounted at the full rate, and lengthens the 4\% rotation to \tbFPriceGrowth{} years; it would reach the
\tbFPriceBoth{} years of the 3\% rotation only if the planting cost rose at 1\% a year as well.

The Octave listing below finds the same rotations without the first order condition, by maximising
the land value directly with \texttt{fminbnd}. It prints the rotations and land values at the three
rates and the rotation with a rising price, and they agree with the figures above to the digits shown.
\Needspace*{19\baselineskip}
\lstinputlisting[language=Octave]{code/ch10/snippets/rotation.m}

\begin{figure}[!tb]
  \centering
  \includegraphics[width=0.86\textwidth]{ch10/fig07_rotation}
  \caption{Land expectation value \cref{eq10:faustmann} of the illustrative stand against the rotation
  length at three discount rates; the dots mark the Faustmann rotations.}
  \label{fig10:rotation}
\end{figure}

\Cref{fig10:marginal} draws the rule of \cref{eq10:foc} for the illustrative stand at $r=4\%$. The value
of a year's growth crosses the interest on the standing timber alone at the single harvest rotation of
\tbSFour{} years. Adding the land rent $r\,\mathit{LEV}$, \mgRent{} dollars a hectare a year, moves the
crossing to the Faustmann rotation of \tbFFour{} years, where a year's growth is worth \mgGrowthValue{}
dollars a hectare.

\begin{figure}[!tb]
  \centering
  \includegraphics[width=0.86\textwidth]{ch10/fig16_marginal}
  \caption{The marginal rule at $r=4\%$ for the illustrative stand: the value of a year's growth $pQ'(T)$
  against the interest on the timber, $rpQ(T)$, with and without the rent on the land at its Faustmann
  value. The crossings are the two rotations.}
  \label{fig10:marginal}
\end{figure}

The value of the land is very sensitive to the discount rate, and that is the main lesson for an
investor. Between 3\% and 6\% the land loses more than nine tenths of its value, because most of its value
lies in harvests decades away. In the model timberland is a long duration asset. In practice the owner can also delay a
harvest, and the value of that flexibility depends on prices, growth, costs and constraints: when timber
prices are low the owner can let the trees grow, and the trees keep growing in volume and
often in value per cubic metre as they reach sizes that sell for more. The option makes timber returns
less tied to the current timber price than the model above, which fixes $p$, suggests.

\subsection{Farmland}

A farm let to a tenant is valued like a building. The rent is set by what the land produces and the prices
of the crops, and \cref{prop10:cap} applies with $g$ the growth of the rent: a farm can give long-horizon
exposure to crop prices, though how much passes through depends on the lease, yields, costs and how the
rent is set. The public data here are the producer price
indexes for farm products and for logs and timber. Measured in December of each year and deflated by the
consumer price index, farm product prices fell by \tbRealFarmAbs\% a year in real terms from \tbFirstFarm{} to
\tbLastYear{} and timber prices by \tbRealTimberAbs\% a year from \tbFirstTimber, with annual real volatilities of
\tbRealVolFarm\% and \tbRealVolTimber\% (\cref{fig10:ppi}). Softwood lumber, a processed product with a more
volatile price, rose by \tbLumberPeak\% from the end of 2019 to its peak in \tbLumberPeakDate. These are prices
of the output and say nothing directly about returns on the land: the owner of farmland or timberland also gains from productivity,
from the growth of the stand and from changes in the land's value, none of which a price index records.

\begin{figure}[!tb]
  \centering
  \includegraphics[width=0.92\textwidth]{ch10/fig08_ppi}
  \caption{Producer price indexes for farm products, for logs, bolts, timber, pulpwood and wood chips,
  and for softwood lumber, each divided by the consumer price index, monthly, log scale.}
  \label{fig10:ppi}
\end{figure}

\subsection{Infrastructure concessions}

A concession gives its holder the cash flows of an asset for $N$ years. With a first year cash flow $X_1$
growing at $g$ and a discount rate $r$, it is a growing annuity.

\begin{proposition}[A concession of finite life]\label{prop10:concession}
For $r>g$ the value of the concession and its Macaulay duration are
\begin{equation}\label{eq10:concession}
  V_N=\frac{X_1}{r-g}\Bigl(1-\Bigl(\frac{1+g}{1+r}\Bigr)^N\Bigr),\qquad
  D_N=\frac{\sum_{t=1}^Nt\,w_t}{\sum_{t=1}^Nw_t},\quad w_t=\frac{(1+g)^{t-1}}{(1+r)^t},
\end{equation}
and as $N\to\infty$, $V_N\to X_1/(r-g)$ and $D_N\to(1+r)/(r-g)$. The share of the perpetuity's value
that the concession captures is $1-\bigl((1+g)/(1+r)\bigr)^N$.
\end{proposition}

\begin{proof}
$V_N=\sum_{t=1}^NX_1w_t$, a geometric sum with first term $X_1/(1+r)$ and ratio $x=(1+g)/(1+r)$, equals
$\frac{X_1}{1+r}\cdot\frac{1-x^N}{1-x}=\frac{X_1}{r-g}(1-x^N)$. For the limit of the duration,
$\sum_{t\ge1}t\,x^{t-1}=1/(1-x)^2$ and $\sum_{t\ge1}x^{t-1}=1/(1-x)$, so $D_\infty=1/(1-x)=(1+r)/(r-g)$.
\end{proof}

With $X_1=\tbCfOne$, $r=\tbRcon\%$ and $g=\tbGcon\%$, the perpetuity is worth \tbPerp. A 30 year concession
is worth \tbConThirty, \tbShareThirty\% of the perpetuity, with a duration of \tbDurThirty{} years; 50 years
gives \tbConFifty{} (\tbShareFifty\%) and 99 years \tbConNinetyNine{} (\tbShareNinetyNine\%), with durations of
\tbDurFifty{} and \tbDurNinetyNine{} years against \tbDurPerp{} for the perpetuity. A rise of one point in the
discount rate lowers the value of the 30 year concession by \tbShockThirtyAbs\% and of the perpetuity by \tbShockPerpAbs\%. When
the tolls are indexed to inflation, $g$ moves with inflation and $r$ is best read as a nominal rate that
moves with it as well, so that the concession is closer to a real bond than to a nominal one. That is the
link to \cref{sec10:inflation}.

\FloatBarrier
% =====================================================================================
\section{Storage and the convenience yield}
\label{sec10:storage}

Chapter~4 priced a forward on a stored commodity by cash and carry: buying the commodity, storing it at a
cost and delivering it, bounds the forward price from above, and the shortfall below that bound defines
the implied convenience yield, the benefit of holding the physical good that \textcite{kaldor1939} named.
Chapter~9 showed that the shape of the futures curve decides the carry earned by a futures position. This
section asks what sets the convenience yield, using two futures prices, which are observed every day,
instead of a spot price, which for crude oil and gas is an average of private trades.

In this section the interest rate $r$ is constant, so futures and forward prices agree, as Chapter~4
showed, and storage costs a share $u$ of the stored quantity per year, paid in kind, so that one unit
stored for $\Delta$ years becomes $e^{-u\Delta}$ units. Payment in kind is a modelling convention; a
storage charge paid in cash changes the carry formula.

\begin{proposition}[Calendar spreads and the price of storage]\label{prop10:storage}
Let $F_1$ and $F_2$ be the futures prices for delivery at $T_1<T_2$, with $\Delta=T_2-T_1$. Suppose storage
is available at the cost $u$, the commodity delivered under either contract can be stored and
redelivered, and the futures can be traded without basis, financing or margin frictions.
\begin{enumerate}[label=(\alph*)]
  \item If there is no arbitrage, $F_2\le F_1e^{(r+u)\Delta}$.
  \item If a holder who has no use for the commodity between $T_1$ and $T_2$ chooses to store it over that
        interval, then $F_2\ge F_1e^{(r+u)\Delta}$, so with (a) the spread is at full carry.
\end{enumerate}
The \emph{net convenience yield} implied by the two contracts, $y-u=r-\ln(F_2/F_1)/\Delta$, is
therefore at least $-u$, so the gross yield $y$ is at least zero. In this frictionless model a strictly
positive implied convenience yield means that no holder without an inventory-service benefit chooses to
carry a marginal unit over this interval.
\end{proposition}

\begin{proof}
(a) Take a long position in one
contract for $T_1$ and a short position in $e^{-u\Delta}$ contracts for $T_2$. At $T_1$ take delivery of
one unit, paying $F_1$ with money borrowed until $T_2$, and store it; at $T_2$ deliver the $e^{-u\Delta}$
units that remain. The strategy costs nothing at the start and pays $e^{-u\Delta}F_2-e^{r\Delta}F_1$ at
$T_2$, which cannot be positive without arbitrage.
(b) A holder of one unit at $T_1$ who stores it has $e^{-u\Delta}$ units at $T_2$. The alternative is to
deliver the unit against a short contract for $T_1$, invest $F_1$ until $T_2$ and buy
$F_1e^{r\Delta}/F_2$ units with a long contract for $T_2$. A holder without other use for the
commodity stores only if storing leaves at least as many units, $e^{-u\Delta}\ge F_1e^{r\Delta}/F_2$,
which is the stated inequality. The last statement rewrites the two inequalities in terms of $y-u$.
\end{proof}

When the curve is inverted, so that the later contract trades below full carry, someone is storing the
commodity at an apparent loss. \textcite{working1949} explained this by the services of inventory: a
refiner or a utility that runs short of stock must stop or buy at any price, so a stock on hand is worth
more to it than a promise of delivery. The marginal value of those services falls as stocks rise, which
gives the \emph{supply of storage} curve of \textcite{working1949} and \textcite{brennan1958}: the price of
storage, $\ln(F_2/F_1)/\Delta-r$, rises with inventories and approaches full carry when stocks are ample.
\textcite{famafrench1987} tested this with the basis of many commodities, and \textcite{gorton2013} relate
the basis and the risk premiums of commodity futures to inventories.

\paragraph{Crude oil}
The US Energy Information Administration published monthly averages of the daily closing prices of the
first four NYMEX crude oil contracts until April 2024, and publishes monthly US crude oil stocks excluding
the Strategic Petroleum Reserve, which here run to \sgStockLast. The first and fourth contracts are about
three months apart, so $\Delta=1/4$, and the 3 month bill rate stands in for $r$. Inventories are
measured against their norm, the log of the stock less the average log stock in the same calendar month
of the previous five years. Averages of daily closes are not prices on a single day, and the delivery
months of the averaged contracts shift within the month, so the measure is approximate.

From \sgOilFirst{} to \sgOilLast, \sgNOil{} months, the net convenience yield had a correlation of
\sgCorrOil{} with the inventory measure (\cref{fig10:working}). In the fifth of months with the lowest stocks,
on average \sgDevLow\% against the norm, it averaged \sgQLow\% a year; in the fifth with the highest
stocks, \sgDevHigh\% above the norm, it averaged \sgQHigh\%. A quadratic fit is convex, steep where stocks
are scarce and flatter where they are ample, the shape Working described, though a fit of this kind
describes the sample and does not by itself identify the schedule of marginal convenience yields. Over the longer record from
\sgBackFirst{} to \sgCurveLast{} the fourth contract was below the first in \sgBackShare\% of months, and the
net convenience yield averaged \sgCyMean\% a year.

\begin{figure}[!b]
  \centering
  \includegraphics[width=\textwidth]{ch10/fig09_working}
  \caption{Net convenience yield implied by the first and fourth NYMEX contracts and the 3 month bill rate
  against inventories measured against the same month of the previous five years, with a quadratic fit.
  (a) WTI crude oil and US crude stocks excluding the SPR, \sgOilFirst{} to \sgOilLast. (b) Henry Hub
  natural gas and US working gas in storage, \sgGasFirst{} to \sgOilLast. Some extreme months lie outside
  the axes.}
  \label{fig10:working}
\end{figure}


April 2020 shows the other end of the curve. As demand collapsed, inventories rose to \sgAprDev\% above
their norm and storage near the delivery point filled; the May contract settled at a negative price on
20 April. In the EIA's monthly averages the first contract averaged \sgAprOne{} dollars a barrel in April
and the fourth \sgAprFour, a net convenience yield of \sgAprCy\% a year. Read mechanically, the frictionless carry bound
of \cref{prop10:storage} with $y\ge0$ would put the annual cost of storage at the margin at \sgAprCyAbs\% or
more. A figure that extreme says the simple model had stopped describing the market: storage near the
delivery point was scarce, and deliverability, location, funding and limits on storage capacity all
entered the spread that month. The spread alone cannot separate them from the cost of storage.

\Cref{fig10:curves} sets April 2020 beside two other months, each curve divided by its first contract,
with dotted lines for the curve that full carry at the bill rate would give with nothing paid for
storage. In April 2020 the fourth contract was \cvAprSpread\% above the first, where the bill would have
carried the price up by \cvAprCarry\%. In \cvBackWhen, the steepest backwardation of any month from January 1990 to \cvLastWhen,
the fourth contract was \cvBackSpreadAbs\% below the first, and in \cvMedWhen, the month with the median
slope, the two were within \cvMedSpread\% of each other.

\begin{figure}[!tb]
  \centering
  \includegraphics[width=0.8\textwidth]{ch10/fig17_curves}
  \caption{Monthly averages of the first four NYMEX crude oil contracts, relative to the first, in three
  months, with the full carry curve at the 3 month bill rate and no storage cost dotted.}
  \label{fig10:curves}
\end{figure}

\FloatBarrier

\paragraph{Natural gas}
Gas is stored in summer and drawn in winter, so its curve has a seasonal shape that crude oil lacks. The
same measure for Henry Hub gas against working gas in storage had a correlation of \sgCorrGas{} over
the \sgNGas{} months with data from \sgGasFirst{} to \sgGasLast; the fourth contract has no average for
\sgGasMissing. The slope from the first to the second contract, annualised, averaged
\sgSlopeSep\% a year in September, when the second contract is the first winter month and the first is
not, and \sgSlopeJan\% in January, when both are late winter months and the second is closer to spring
(\cref{fig10:gas}). Part (a) of \cref{prop10:storage} limits how steep the curve can be from summer into
winter, through the cost of storage, but the proposition gives no lower bound on how far it can fall from winter into
spring, because the reverse trade would need gas that can be borrowed and returned later, and gas held
for winter use is not lent out that way.

\begin{figure}[!tb]
  \centering
  \includegraphics[width=0.82\textwidth]{ch10/fig10_gas}
  \caption{Annualised log slope from the first to the second Henry Hub natural gas contract, averaged by
  calendar month, \sgGasCurveFirst{} to \sgCurveLast.}
  \label{fig10:gas}
\end{figure}

\FloatBarrier

\subsection{The value of a storage facility}
\label{sec10:facility}

A gas storage field or an oil tank earns the difference between the price at which its owner fills it
and the price at which the owner empties it. The futures curve lets the owner fix much of that difference
in advance, and the rest of the facility's value comes from changing the plan as the curve moves.

\begin{proposition}[Intrinsic value of storage]\label{prop10:intrinsic}
A facility holds at most one unit, costs $c_{\mathrm{in}}$ per unit to fill and $c_{\mathrm{out}}$ to
empty, is empty at time $0$ and has no other cost. Filling and emptying are instantaneous, nothing is lost in
storage, the commodity delivered under the contracts can be stored in the facility, and the futures can
be traded at the quoted prices without costs or margin constraints. Futures prices $F_1,\dots,F_n$ for delivery at
$t_1<\dots<t_n$ are quoted at time $0$, and $r$ is constant.
\begin{enumerate}[label=(\alph*)]
  \item At no cost at time $0$, the owner can fix one cycle of filling and emptying worth the single-cycle \emph{intrinsic value}
        \begin{equation}\label{eq10:intrinsic}
          I=\max\Bigl(0,\;\max_{i<j}\bigl[e^{-rt_j}(F_j-c_{\mathrm{out}})-e^{-rt_i}(F_i+c_{\mathrm{in}})\bigr]\Bigr).
        \end{equation}
  \item If the owner may later close futures positions and open new ones at the prices then quoted, the
        value of running the facility with such revisions is at least $I$, since the owner may always keep
        the plan already fixed, and revisions add nothing when futures prices never change.
\end{enumerate}
\end{proposition}

\begin{proof}
(a) Buy one contract for $t_i$ and sell one for $t_j$, for the pair that attains the inner maximum.
Futures cost nothing to open, and with $r$ constant they are worth the same as forwards, as Chapter~4
showed, so the position is a forward purchase at $F_i$ and a forward sale at $F_j$. At $t_i$ the owner
takes delivery, pays $F_i$ and the cost of filling and stores the unit; at $t_j$ the owner pays the cost of
emptying and delivers it for $F_j$. The present value at $0$ is the bracket. If no bracket is positive,
the owner leaves the facility empty and $I=0$.
(b) At a later date $s$ the owner holds the futures positions of a plan fixed earlier and some
inventory. Compare plans by the discounted value at $s$ of their self-financing cash flows, keeping the
variation margin already accrued: that margin, paid or received up to $s$ and carried at $r$, is the same
whatever the owner does next, so it enters every comparison equally. After the settlement at $s$ an open
position is a claim to its remaining margin and its delivery payment, which with $r$ constant have the
present value of a delivery at the futures price of date $s$. Changing a hedge alters those remaining
cash flows, but a position opened or closed at that price has zero value. The rest of the current plan is
therefore worth the present value, at the prices of date $s$, of its remaining flows, which is its value
as one of the plans open to the owner at $s$, given what is in the facility. Switching to the best plan at
those prices costs nothing and replaces that value by the value of the best plan, which is never smaller,
since the current plan is among those open. Every revision therefore adds a discounted gain that is zero or positive to the $I$ fixed at
time $0$. If prices never change, the best plan at every date is the one already held and every gain is
zero.
\end{proof}

The excess of the facility's value over $I$ is called its extrinsic value. It comes from the movement of
the curve; it generally benefits from greater spread volatility when the price process and operational
constraints are otherwise held fixed, but its value requires an explicit stochastic model. With several cycles,
linear costs and linear limits on how fast the facility can be filled or emptied, the inner maximum of
\cref{eq10:intrinsic} becomes a linear programme over the quantities injected and withdrawn each month, and
the argument for (b) is unchanged; valuing the revisions themselves needs a model of how the curve moves.

The listing below sets up that linear programme for one unit of capacity, with injections and
withdrawals at each delivery date and the inventory kept between empty and full. On the curve of
\cref{exr10:fifteen} it recovers the intrinsic value of \xoV{} from the best pair of months. On a
stylised seasonal curve of twelve monthly deliveries, peaking in the fifth, it prints \xoSeasOne{} when
the facility can be filled in a month and \xoSeasThird{} when it can be filled or emptied by only a third
of its capacity a month: the rate limit removes almost half the value.
\Needspace*{26\baselineskip}
\lstinputlisting[language=Python]{code/ch10/snippets/storage_lp.py}

The EIA averages give a test of part (a) for gas. In each August from \sgLockFirst{} to \sgLockLast,
the first contract was, for most of the month, the September contract and the fourth the December one,
and their monthly averages describe an indicative September to December spread. An owner could have
fixed an actual spread close to it by trading the pair on a day of that month, at that day's prices and
within the facility's limits. The log spread of the averages was \sgLockMean\% on average, with a standard
deviation of \sgLockSd\%, never fell below \sgLockMin\%. An owner who instead filled in September and emptied in December at the monthly
average spot price earned a spread that averaged \sgRealMean\% with a standard deviation of \sgRealSd\%,
and was positive in \sgRealPos\% of years (\cref{fig10:lock}). The spread fixed in August is known when the
facility is filled; the spot spread depends on the weather of the coming winter. The fixed spread was on
average \sgGapMean{} percentage points above the spot spread that followed, with a standard error of
\sgGapSe{} points, a $t$ statistic of about \sgGapT. That agrees with a premium paid by those who buy
winter gas forward to those who store it, though \sgLockN{} years of one market is thin evidence for its
size. Over a whole storage season, July to January, the
spot spread averaged \sgJanMean\% over the \sgJanN{} seasons from \sgJanFirst{} to \sgJanLast, with a
standard deviation of \sgJanSd\%.

\begin{figure}[!tb]
  \centering
  \includegraphics[width=\textwidth]{ch10/fig12_lock}
  \caption{Henry Hub natural gas: the log spread from the first to the fourth NYMEX contract in the August
  average, about September to December, and the change in the monthly average spot price from September to
  December of the same year, \sgLockFirst{} to \sgLockLast.}
  \label{fig10:lock}
\end{figure}

\FloatBarrier
% =====================================================================================
\section{Inflation hedging}
\label{sec10:inflation}

Real assets are often bought as a hedge against inflation, on the view that rents, crop prices and the
prices of stored goods rise with the general price level. A hedge in this sense is an asset whose real
return does not fall when inflation rises. The simplest way to measure it is a regression of the asset's
nominal return on inflation over the same period,
\begin{equation}\label{eq10:beta}
  R_t = a + b\,\pi_t + e_t, \qquad \Cov(e_t,\pi_t)=0,
\end{equation}
and \textcite{fama1977} refined it by splitting inflation into the part expected at the start of the period
and the surprise.

\begin{proposition}[Inflation beta and the real return]\label{prop10:hedge}
Under \cref{eq10:beta}, the real return $R_t-\pi_t$ satisfies
\[
  \Cov(R_t-\pi_t,\pi_t)=(b-1)\Var\pi_t,\qquad
  \Var(R_t-\pi_t)=(b-1)^2\Var\pi_t+\Var e_t .
\]
The real return is uncorrelated with inflation if and only if $b=1$, and the asset's real return varies
least for a given $\Var e_t$ when $b=1$.
\end{proposition}

\begin{proof}
$R_t-\pi_t=a+(b-1)\pi_t+e_t$, with $e_t$ uncorrelated with $\pi_t$; both formulas follow, and for a fixed
$\Var(e_t)$ the variance is smallest at $b=1$.
\end{proof}

A beta of one makes an asset a hedge in the sense that its real return is uncorrelated with inflation
under the regression model; independence would need more assumptions about the residual, and it does not make the asset safe, because the residual $e_t$ remains. A nominal bill held for a year has a
beta close to one on expected inflation, since its yield is set with inflation in view, and close to zero
on the surprise. The question for an investor who holds bills is how much of a real asset to add to make
the real return of the mix as steady as possible.

\begin{proposition}[The least real variance mix]\label{prop10:mix}
Let $f_t$ be the return on a bill known at the start of period $t$, $x_t$ the return on a risky asset and
$\pi_t$ inflation. The weight $w$ in the risky asset, with $1-w$ in the bill, that minimises the variance of
the real return $(1-w)f_t+wx_t-\pi_t$, with $\Var(x_t-f_t)>0$, is
\begin{equation}\label{eq10:mix}
  w^*=\frac{\Cov(x_t-f_t,\;\pi_t-f_t)}{\Var(x_t-f_t)},
\end{equation}
and it reduces the variance of the real return of the bill alone by the share $\varrho^2$, where $\varrho$ is
the correlation of $x_t-f_t$ with $\pi_t-f_t$. The weight is unconstrained and may be negative or above
one; if short positions or borrowing are ruled out, the best permitted weight is $w^*$ moved to the
nearest end of the allowed range.
\end{proposition}

\begin{proof}
The linearised real return is $w(x_t-f_t)-(\pi_t-f_t)$, the error left when $w$ units of $x_t-f_t$
hedge $\pi_t-f_t$; at $w^*$ it is, up to a constant, minus the least squares residual of $\pi_t-f_t$ on
$x_t-f_t$. Its variance $\Var(\pi-f)-2w\Cov(x-f,\pi-f)+w^2\Var(x-f)$ is a quadratic in $w$ minimised at $w^*$, where it
equals $\Var(\pi-f)(1-\varrho^2)$.
\end{proof}

The objective is linear in $w$ by construction. With simple returns the mix earns exactly
$(1-w)f_t+wx_t$, but subtracting $\pi_t$ gives the real return only to first order; with log rates the
subtraction is exact, but the weighted sum is only a first-order log return of the rebalanced mix.
Either way \cref{eq10:mix} is exact for the linear objective and first order for the mix.

\paragraph{The evidence}
\Cref{tab10:inflation} reports \cref{eq10:beta} for calendar years from \ihFirst{} to \ihLast, or from the
first full year of each series, with Newey and West standard errors over two lags. Inflation is the
December to December change in the CPI, which averaged \ihInflMean\% a year with a standard deviation of
\ihInflSd\%. The bill is a 3 month Treasury bill rolled each month. Property, houses, farm products,
timber and oil are price changes without income, so their mean real returns understate what an owner
earned; equities and bills are total returns. Following \textcite{fama1977}, the bill yield at the end of the previous December
serves as a proxy for expected inflation, though it also holds expected real interest and other
components, and the surprise is the rest, so it too is a proxy and not an observed forecast error.

The R listing below reproduces the property row of \cref{tab10:inflation} from the data files, with
the Newey and West covariance written out: Bartlett weights over two lags and no small sample factor,
as in the Python code. It prints the \ihNCre{} years, the slope of \ihBCre{} and its standard error of
\ihSeCre.
\Needspace*{20\baselineskip}
\lstinputlisting[language=R]{code/ch10/snippets/inflation.R}

\begin{table}[t]
  \centering
  \small
  \caption{Calendar year log returns and price changes against CPI inflation, \ihFirst{} to \ihLast{} or from
  the year shown. Mean and volatility of the real return in percent; the slope on inflation alone, and
  the slopes on expected and unexpected inflation from one regression using the bill-yield proxy, with Newey and West standard errors in brackets; cumulative real
  return in 2021 and 2022, when inflation was \ihInflTwo\%. Price changes exclude income.}
  \label{tab10:inflation}
  \setlength{\tabcolsep}{3pt}
  \begin{tabular}{@{}lrrrrrrr@{}}
    \toprule
    & From & Real mean & Real vol. & Slope on $\pi$ & On expected & On surprise & 2021--22\\
    \midrule
    3 month bills & \ihStartBill & \ihRealBill & \ihVolBill & \ihBBill{} (\ihSeBill) & \ihEBill{} (\ihEseBill) & \ihUBill{} (\ihUseBill) & \ihTwoBill\\
    Commercial property & \ihStartCre & \ihRealCre & \ihVolCre & \ihBCre{} (\ihSeCre) & \ihECre{} (\ihEseCre) & \ihUCre{} (\ihUseCre) & \ihTwoCre\\
    Houses & \ihStartHpi & \ihRealHpi & \ihVolHpi & \ihBHpi{} (\ihSeHpi) & \ihEHpi{} (\ihEseHpi) & \ihUHpi{} (\ihUseHpi) & \ihTwoHpi\\
    Equity REITs, price & \ihStartReit & \ihRealReit & \ihVolReit & \ihBReit{} (\ihSeReit) & \ihEReit{} (\ihEseReit) & \ihUReit{} (\ihUseReit) & \ihTwoReit\\
    US equities & \ihStartMkt & \ihRealMkt & \ihVolMkt & \ihBMkt{} (\ihSeMkt) & \ihEMkt{} (\ihEseMkt) & \ihUMkt{} (\ihUseMkt) & \ihTwoMkt\\
    Farm products & \ihStartFarm & \ihRealFarm & \ihVolFarm & \ihBFarm{} (\ihSeFarm) & \ihEFarm{} (\ihEseFarm) & \ihUFarm{} (\ihUseFarm) & \ihTwoFarm\\
    Logs and timber & \ihStartTimber & \ihRealTimber & \ihVolTimber & \ihBTimber{} (\ihSeTimber) & \ihETimber{} (\ihEseTimber) & \ihUTimber{} (\ihUseTimber) & \ihTwoTimber\\
    Crude oil & \ihStartOil & \ihRealOil & \ihVolOil & \ihBOil{} (\ihSeOil) & \ihEOil{} (\ihEseOil) & \ihUOil{} (\ihUseOil) & \ihTwoOil\\
    \bottomrule
  \end{tabular}
\end{table}

Bills followed expected inflation, with a slope of \ihEBill{} on it, and missed the surprises, with a slope
of \ihUBill; in 2021 and 2022 they lost \ihTwoBillAbs\% in real terms. Commercial property and house prices had
slopes on inflation of \ihBCre{} and \ihBHpi, and slopes close to one on the surprise proxy, \ihUCre{} and \ihUHpi, over this sample, as \textcite{fama1977} found for private
residential real estate, though with standard errors of about
0.3 the data cannot tell a slope of one from a slope of a half. Farm products and timber had slopes of
\ihBFarm{} and \ihBTimber, and farm prices rose with the surprises by more than one for one. Oil's slope of
\ihBOil{} is partly mechanical, because energy prices are part of the CPI, and a regression of this kind
cannot say which way the influence runs, though energy's weight in the CPI means some of it must run from oil to the index. US equities had a negative slope of \ihBMkt, as
\textcite{bodie1976} and \textcite{fama1977} found in earlier data, with a standard error too large to say
much. The REIT sample is short and its slope is not estimated with any precision.

\Cref{fig10:betas} draws the slopes with two standard errors. Taken together, several of these price series have estimated slopes on
inflation that are positive and often near one, though imprecise and without owner income, and the
hedge is noisy: the residual volatility is large beside the inflation it is meant to offset. For
property, \cref{prop10:mix} gives a weight of \ihWcre{} in the property price index, with the rest in
bills, as the mix with the least real variance, and it cuts the variance of the bills' real return, a
volatility of \ihRealVolBill\%, by only \ihVarCut\%. A real asset is a hedge against inflation in the sense
of \cref{prop10:hedge} and a poor one in the sense of \cref{prop10:mix}, because the risks that come with
it are much larger than the inflation risk it removes.

\begin{figure}[!tb]
  \centering
  \includegraphics[width=0.86\textwidth]{ch10/fig11_betas}
  \caption{Slopes of calendar year returns and price changes on CPI inflation, from the year shown to
  \ihLast, with two Newey and West standard errors.}
  \label{fig10:betas}
\end{figure}

\Cref{fig10:crescatter} shows why for property. The fitted slope of \ihBCre{} lies close to the one for
one line, but the years scatter around it with a residual standard deviation of \scResSd{} points,
three times the \ihInflSd{} point standard deviation of inflation itself. The largest miss, in \scBigYear, was \scBigRes{} points.

\begin{figure}[!htb]
  \centering
  \includegraphics[width=0.66\textwidth]{ch10/fig18_crescatter}
  \caption{Calendar year log changes in the Z.1 commercial property price index against December to
  December CPI inflation, \ihFirst{} to \ihLast, with the fitted line of \cref{tab10:inflation} and a one for
  one line.}
  \label{fig10:crescatter}
\end{figure}

\FloatBarrier

\subsection{Longer horizons}
\label{sec10:horizon}

An investor who holds a building for ten years cares about its real value at the end more than about
each year's return. The slope on inflation can differ between horizons, and there is a simple reason to
expect it to settle for an asset whose real price returns to a trend.

\begin{proposition}[Inflation slopes at long horizons]\label{prop10:horizon}
Let $p_t$ be the log price level and $v_t$ the log price of an asset, with $v_t=\theta p_t+\delta t+u_t$,
where $u_t$ is stationary with variance $\sigma_u^2$, and suppose that $\Var(p_{t+k}-p_t)\to\infty$ as
$k\to\infty$. Write $\Delta_k x_t=x_t-x_{t-k}$. The slope of the $k$ period change in $v$ on the $k$
period change in $p$,
\[
  b_k=\frac{\Cov(\Delta_k v_t,\Delta_k p_t)}{\Var(\Delta_k p_t)},
  \qquad\text{satisfies}\qquad
  \lvert b_k-\theta\rvert\le\frac{2\sigma_u}{\sqrt{\Var(\Delta_k p_t)}},
\]
so $b_k\to\theta$. If $u_t$ is uncorrelated with the price level at all leads and lags, $b_k=\theta$ at
every horizon.
\end{proposition}

\begin{proof}
$\Delta_k v_t=\theta\Delta_k p_t+\delta k+\Delta_k u_t$, and the constant $\delta k$ has no covariance with
anything, so $b_k-\theta=\Cov(\Delta_k u_t,\Delta_k p_t)/\Var(\Delta_k p_t)$. By the Cauchy and Schwarz
inequality the covariance is at most $\sqrt{\Var(\Delta_k u_t)\Var(\Delta_k p_t)}$ in absolute value, and
$\Var(\Delta_k u_t)=2\sigma_u^2(1-\rho_k)\le4\sigma_u^2$, where $\rho_k$ is the $k$th autocorrelation of
$u_t$. Dividing gives the bound, which goes to zero because $\Var(\Delta_k p_t)$ grows without limit. If
$u_t$ is uncorrelated with $p_s$ for all $s$, the covariance is zero for every $k$.
\end{proof}

The condition on the price level holds when inflation is stationary with a long run variance above zero,
for then $\Var(\Delta_k p_t)$ grows in proportion to $k$. The slope moves with the horizon only through the
covariance of the gap $u_t$ with inflation: an asset whose price follows the price level with a lag, as
the rent on a lease fixed in nominal terms does, has a gap that is low after a burst of inflation, so its
slope is below $\theta$ over short horizons and approaches $\theta$ over long ones (Exercise~\ref{exr10:sixteen}).
The bound is a worst case and is slow: it shrinks only with the square root of the horizon.
\textcite{boudoukh1993} found that the relation between stock returns and inflation, negative over short
horizons, turns positive over five years in two centuries of US and UK data.

\Cref{fig10:horizon} repeats the regressions of \cref{tab10:inflation} with overlapping sums over three,
five and ten years from \ihFirst{} to \ihLast, with Newey and West standard errors over as many lags as the
horizon. With fifty five years, neighbouring ten year sums have nine of their ten years in common, so the \ihHnTen{}
overlapping observations carry far less information than \ihHnTen{} independent ones, and the standard
errors are themselves imprecise. Bills,
rolled at rates that follow expected inflation, had slopes of \ihHBillThree, \ihHBillFive{} and \ihHBillTen.
US equities moved as the lag story suggests, from \ihBMkt{} over a year to \ihHMktFive{} over five years and
\ihHMktTen{} over ten, with a standard error of \ihHsMktTen{} at ten. The property index did not: its slope
was \ihHCreThree, \ihHCreFive{} and \ihHCreTen{} at three, five and ten years, with standard errors of about
\ihHsCreTen{} to \ihHsCreThree, and house prices had slopes of \ihHHpiThree, \ihHHpiFive{} and \ihHHpiTen. In the terms of the proposition that
points to $\theta$ below one for property prices, which the standard errors cannot separate from one.
Farm product prices had smaller slopes at longer horizons, \ihHFarmTen{} at ten years, and timber larger
ones, \ihHTimberTen{} with a standard error of \ihHsTimberTen. The REIT series, which starts in 1997 when
inflation varied little, gives slopes too imprecise to report.

\begin{figure}[!tb]
  \centering
  \includegraphics[width=0.86\textwidth]{ch10/fig13_horizon}
  \caption{Slopes of overlapping one, three, five and ten year sums of returns and price changes on the
  sums of CPI inflation over the same years, \ihFirst{} to \ihLast, or from the first year of each series.}
  \label{fig10:horizon}
\end{figure}

A simulation shows how far those standard errors can be trusted. In the model of
Exercise~\ref{exr10:sixteen} an asset's log price change is last year's inflation, so the true slope
over $k$ years is $1-1/k$. Draw inflation as independent normals with the mean and standard deviation of
\ihFirst{} to \ihLast, and take \lsN{} samples of \lsT{} years from the chapter's random stream
\texttt{[3817, 3]}. The simulated ten year slope averaged \lsMeanTen{} against a true 0.9, with a standard
deviation across samples of \lsSdTen; the median Newey and West standard error was \lsSeTen, and a test at
the 5\% level rejected the true slope in \lsRejTen\% of samples, against \lsRejFive\% at five years and
\lsRejOne\% at one. These come from an assumed model and say nothing directly about the data, but they show that,
under this simulated lagged-price model, the specified Newey and West procedure understates ten year
sampling uncertainty, so a ten year slope two standard errors away from one is weaker evidence than it
looks; this does not establish the direction of the error in the empirical regressions.

\FloatBarrier
% =====================================================================================
\section{Limitations}
\label{sec10:limits}

The building, its loan and the forest are illustrations with round inputs, and the simulation of
\cref{sec10:leverage} draws rent growth and exit cap rates from normal laws with no link between them; in a
real downturn rents fall and cap rates rise together, and the tails of the levered IRR would be wider. The
loan case study marks single loans to a national index, which moves far less than a single building, uses
one spread for twenty five years and ignores the covenants and extensions that lenders use when values fall.

The property series are price indexes. The Z.1 index and the FHFA index measure prices of buildings
and houses, with no rent and no allowance for the capital spending needed to keep a building's value,
and the REIT series is a revaluation of market value without dividends. Total returns would be higher by
the income yield, and their correlations with inflation and with the market could differ. The splice of
the Z.1 index in 1996 changes the method and the property mix at the same time, so not all of the change in
autocorrelation can be attributed to appraisals. The leverage used to unlever REITs is an assumption
over a range, and REIT leverage has varied over time.

The storage results rest on monthly averages of daily futures prices that the EIA stopped publishing in
April 2024, on national inventories that mix grades and locations, and on a bill rate that is not the
financing rate of a commodity trader. The fixed and realised gas spreads use \sgLockN{} years of monthly averages, and the realised spread
is measured at spot prices that a storage owner would not obtain exactly. The inflation regressions use
about fifty annual observations at most, and fewer for REITs and oil, so their standard errors are large; the inflation of 2021 and 2022 is a
single episode, and the overlapping sums over long horizons add little independent information. The producer price indexes are prices of output, not returns on owning farmland or
timberland, and the rotation model fixes the price of timber that in practice moves and can be waited out.

\FloatBarrier
% =====================================================================================
\section{Summary}
\label{sec10:summary}

\begin{itemize}
  \item The cap rate is next year's NOI over value. With constant growth and a capital reserve it equals
        $(r-g)/(1-\kappa)$, and a buyer who pays that price and sells at the same cap rate earns
        $(1-\kappa)c+g$, whatever the holding period.
  \item In a discounted cash flow the exit cap rate carries a share $((1+g)/(1+r))^T$ of the value, which
        for ten years at \pfR\% and \pfG\% is \pfSaleShare\%.
  \item A lender's DSCR test is $c/(\lambda k)\ge\underline D$. When the mortgage constant is high against
        the cap rate, as for the building's loan at \lvRate\%, whose constant of \lvK\% is above the \pfCap\% cap rate, the DSCR test sets the loan, here at \lvLtv\% of value.
        Leverage at a cost below the expected property return raises the expected equity return, multiplies
        its volatility by $1/(1-\lambda)$ and can still lower the first year's cash yield.
  \item Marked to the Z.1 index, \lcShareEighty\% of loans written at \lcLtvZero\% from 1996 passed 80\%
        during their five year term; the loan of \lcWorst{} reached \lcWorstPeak\%.
  \item In the smoothing model, appraisal smoothing changes the cumulative return only through its
        end points and understates volatility, beta and correlation. The Z.1 property index had a first autocorrelation of \smArPre{} and a
        volatility of \smVolPre\% when it was built from appraisals, and \smArPost{} and \smVolPost\% after it
        moved to repeat sales.
  \item Equity REIT returns are most correlated with the property index's later changes. Their
        correlation with the index rises from \rtCorrCreReit{} in the same quarter to \rtLeadOne{} against
        the previous quarter's REIT return, and in annual data to \rtCorrYLag{} against the previous
        year's.
  \item The Faustmann rotation equates the value of a year's growth with the interest on the timber and
        the rent on the land, and is shorter than the rotation for a single harvest. A finite concession of
        30 years at \tbRcon\% and \tbGcon\% growth is worth \tbShareThirty\% of the perpetuity.
  \item Futures spreads are bounded by full carry. An inverted curve with stocks in storage implies a
        convenience yield, and the net convenience yield of crude oil fell as inventories rose, with a
        correlation of \sgCorrOil. The owner of storage can lock in an available spread with futures and keeps the option
        to revise the plan if the curve moves in its favour, an option whose value depends on how prices move
        and on the facility's limits; the August monthly average gas futures imply a September to
        December log spread of \sgLockMean\% on average, against a spot spread of \sgRealMean\% with a volatility of \sgRealSd\%.
  \item Property, house, farm and timber prices had inflation slopes between \ihBHpi{} and \ihBFarm{} in
        annual data from the 1970s, or \ihStartTimber{} for timber, and equities a negative slope that is not significant, but the residual risk is large: the property mix with least real variance holds
        \ihWcre{} in property and cuts the real variance of bills by \ihVarCut\%.
  \item Over long horizons the slope of an asset whose real price returns to a trend tends to the
        weight $\theta$ of the price level in its price. Equity slopes rose from \ihBMkt{} over a year to
        \ihHMktTen{} over ten; property slopes fell from \ihBCre{} to about \ihHCreTen.
\end{itemize}

\FloatBarrier
% =====================================================================================
\section{Exercises}
\label{sec10:exercises}

\begin{enumerate}[label=\textbf{10.\arabic*}, ref=10.\arabic*, leftmargin=*]
  \item\label{exr10:one} A building has a potential rent of 1.2 million dollars, vacancy of 5\% and
  operating expenses of 0.40 million. Its owner reserves 8\% of NOI for capital items, requires 8\% and
  expects NOI to grow at 2\% a year. Find NOI, the cap rate and the value.
  \item\label{exr10:two} A ten year DCF uses $r=7\%$, $g=2\%$ and a going-in cap rate of 5\% with no
  reserve. What share of the value comes from the sale? By how much does the value change if the exit cap
  rate is 5.5\%, and how close is the approximation of \cref{prop10:exit}?
  \item\label{exr10:three} A property is bought at a cap rate of 6\% with NOI growing at 2\% and no
  reserve. What is the IRR if it is sold at the same cap rate? What is it if it is sold after seven years at
  6.5\%?
  \item\label{exr10:four} Find the mortgage constant of a loan at 5.5\% amortised monthly over 30 years,
  and the DSCR of a loan at 70\% of value on a property with a cap rate of 6\%.
  \item\label{exr10:five} A property earns NOI of 2 million dollars and is worth 32 million. A lender
  offers 6\% over 25 years and requires a DSCR of at least 1.30 and an LTV of at most 70\%. Which test
  binds, and how large is the loan?
  \item\label{exr10:six} What share of a 25 year loan at 6\% is still owed after ten years?
  \item\label{exr10:seven} A property returns 7\% and debt costs 5.5\%. Find the equity return at 60\%
  LTV and the equity volatility if the property's volatility is 10\%. What is the equity return if the
  property returns 4\%?
  \item\label{exr10:eight} The true quarterly returns of a property portfolio are independent with an
  annual volatility of 10\%, and an appraisal index smooths them with $\alpha=0.7$. What annual volatility
  does the index report, scaling its quarterly volatility by $\sqrt4$, and what is its fourth order
  autocorrelation?
  \item\label{exr10:nine} An index smoothed with $\alpha=0.6$ reports quarterly returns of 1.0, 1.5, 0.5
  and $-1.0$ percent. Recover the true returns of the last three quarters.
  \item\label{exr10:ten} Smoothing with $\alpha=0.5$ is applied to a property whose true beta is one.
  What is the sum of the slopes in a Dimson regression with the current market return and three lags?
  \item\label{exr10:eleven} For the forest of \cref{sec10:timber}, find the Faustmann rotation, the single
  harvest rotation and the land value at $r=5\%$.
  \item\label{exr10:twelve} A 25 year concession pays 50 in its first year, growing at 3\%, and is discounted
  at 8\%. Find its value, the share of the perpetuity's value and its duration.
  \item\label{exr10:thirteen} Crude oil trades at 70 dollars for delivery in one month and 71 for delivery
  six months after that. The interest rate is 4\% and storage costs 2\% a year. What is the highest
  price for the later contract that rules out arbitrage, and what net and gross convenience yields do the
  prices imply?
  \item\label{exr10:fourteen} An asset has an inflation beta of 0.6 and inflation has a standard deviation
  of 3\%. Find the covariance of its real return with inflation. If the asset's return in excess of a bill
  has a volatility of 10\% and a covariance of $0.6\times0.03^2$ with inflation in excess of the bill, what
  weight in the asset minimises the variance of the real return?
  \item\label{exr10:fifteen} In September, gas futures for delivery in October, January and April trade at
  3.00, 3.60 and 3.20 dollars, for delivery in one, four and seven months. The interest rate is 4\%, and
  filling and emptying a facility of one unit each cost 0.05. The facility starts and ends empty,
  filling and emptying are instantaneous, and the quoted prices can be traded at once. Find the intrinsic value of one cycle and the
  pair of months that attains it. What would a cycle from October to April be worth?
  \item\label{exr10:sixteen} Inflation $\pi_t$ is independent from year to year with variance
  $\sigma^2$, and an asset's log price is last year's log price level, $v_t=p_{t-1}$. Find $\theta$ and
  $u_t$ in \cref{prop10:horizon}, and the slope $b_k$ for $k=1$, $5$ and $10$ years.
\end{enumerate}

\subsection*{Answers}
\begin{enumerate}[label=\textbf{10.\arabic*}, leftmargin=*]
  \item Effective gross income is $1.2\times0.95=1.14$ million and NOI is \xaNoi{} thousand. By
  \cref{prop10:cap} the cap rate is $0.06/0.92=\xaCap\%$ and the value $0.74\times0.92/0.06=\xaValue$ million.
  \item The share is $(1.02/1.07)^{10}=\xbShare\%$. At an exit cap of 5.5\% the value changes by \xbDrop\%,
  and the elasticity gives a change in log value of $-\xbShare\%\times\ln1.1=\xbApprox\%$, which is a
  change of \xbApproxPct\% in value, close because the change is small.
  \item By \cref{prop10:dcf} the IRR is $6+2=\xcIrr\%$. With the sale at 6.5\% after seven years it is \xcIrrHigh\%.
  \item The constant is \xdK\%, and the DSCR is $0.06/(0.70\times\xdK\%)=\xdDscr$.
  \item The constant is \xeK\%. The DSCR test allows $2/(1.30\times\xeK\%)=\xeLoan$ million and the LTV test
  \xeLtvLoan{} million, so the DSCR test binds and the loan is \xeLtv\% of value.
  \item By \cref{eq10:balance}, \xfBal\% of the loan.
  \item $(7-0.6\times5.5)/0.4=\xgRet\%$, with volatility $10/0.4=\xgEqVol\%$. At 4\% the equity earns
  $(4-3.3)/0.4=\xgRetLow\%$, less than the property: leverage is negative.
  \item The reported variance is $0.3/1.7$ of the true one, so the reported volatility is
  $10\%\times\sqrt{0.3/1.7}=\xhVol\%$, and the fourth order autocorrelation is $0.7^4=\xhAcf$.
  \item By \cref{eq10:unsmooth}, $(1.5-0.6)/0.4=\xiOne$, $(0.5-0.9)/0.4=\xiTwo$ and
  $(-1.0-0.3)/0.4=\xiThree$ percent. The true returns swing far more than the reported ones.
  \item The slopes are $0.5\times0.5^j$, and their sum over $j=0,\dots,3$ is \xjSum.
  \item The Faustmann rotation is \xkF{} years, the single harvest rotation \xkS{} years, and the land is
  worth \xkLev{} dollars a hectare.
  \item By \cref{prop10:concession} the value is \xlV, \xlShare\% of the perpetuity's $50/0.05=1000$, and the
  duration is \xlDur{} years.
  \item The bound of \cref{prop10:storage} is $70e^{0.06\times0.5}=\xmBound$. The net convenience yield is
  $0.04-\ln(71/70)/0.5=\xmCy\%$ a year, and adding back the storage cost gives a gross yield of \xmGross\%.
  \item By \cref{prop10:hedge} the covariance is $(0.6-1)\times0.03^2=\xnCov$. By \cref{prop10:mix} the weight
  is $0.6\times0.0009/0.01=\xnW$.
  \item By \cref{eq10:intrinsic}, October to January is worth
  \[e^{-0.04/3}(3.60-0.05)-e^{-0.04/12}(3.00+0.05)=\xoV,\]
  the largest of the three pairs. October to April is worth only \xoLast, and January to April is worth
  nothing, since the curve falls into spring. This is the discounted payoff of the best single cycle under
  the exercise's assumptions; a full valuation of a storage facility would add its extrinsic value.
  \item $v_t=p_t-\pi_t$, so $\theta=1$, $\delta=0$ and $u_t=-\pi_t$. Then $\Delta_kv_t=\Delta_kp_t-\pi_t+\pi_{t-k}$.
  The sum $\Delta_kp_t$ holds $\pi_t$ but not $\pi_{t-k}$, so the covariance is $k\sigma^2-\sigma^2$ and
  $b_k=1-1/k$: zero over a year, 0.8 over five and 0.9 over ten.
\end{enumerate}

\FloatBarrier
% =====================================================================================
\section{Notes, sources and credits}
\label{sec10:credits}

\paragraph{Sources}
The growing perpetuity behind the cap rate is from \textcite{gordon1956}. The smoothing model and its
inversion are from \textcite{geltner1991,geltner1993}, the lagged regression for betas from
\textcite{dimson1979}, and the smoothing of hedge fund returns from \textcite{getmansky2004}. The comparison
of public and private real estate after restating for leverage, property mix and smoothing is from
\textcite{pagliari2005}. The qualification rules for US REITs are in sections 856 and 857 of the Internal
Revenue Code \parencite{usc856}. The forest rotation is from \textcite{faustmann1995}, an English translation
of the article of 1849, with the history and the economics in \textcite{samuelson1976}; the growth curve is
from \textcite{richards1959}. Convenience yield was named by \textcite{kaldor1939}, the theory of storage is
from \textcite{working1949} and \textcite{brennan1958}, and its tests on futures bases are in
\textcite{famafrench1987} and, with inventories, \textcite{gorton2013}. Commodity futures as an asset class are
discussed by \textcite{gorton2006} and \textcite{erb2006}. Inflation betas and the split into expected and
unexpected inflation are from \textcite{fama1977}, and the negative relation of stock returns to inflation
in \textcite{bodie1976}, with its reversal over long horizons in \textcite{boudoukh1993}. The cap rate and exit cap results, the loan tests, the leverage and unlevering
formulas, the smoothing identity, the concession formulas, the calendar spread bounds, the intrinsic value of storage, the least real
variance mix and the long horizon bound are short derivations given here for completeness.

\paragraph{Software used}
The scripts call NumPy \parencite{numpy} for arrays and random numbers, SciPy \parencite{scipy} for root
finding and optimisation, pandas \parencite{pandas} to read and align the series and to parse the EIA
tables, statsmodels \parencite{statsmodels} for the Newey and West standard errors, and Matplotlib
\parencite{matplotlib} for the figures. The cross-checks run in R \parencite{rcore} and GNU Octave
\parencite{octave}. \texttt{code/requirements-lock.txt} records the full environment.

\paragraph{Projects for comparison}
The numpy-financial package \parencite{numpyfinancial} and pyxirr \parencite{pyxirr} compute the IRRs and
loan payments of \cref{sec10:value,sec10:leverage}, and can be used to check the code here. The full list, with
links, is in \texttt{CREDITS.md}.

\paragraph{Data}
The commercial property price index and the equity REIT market value and revaluation are from the
Financial Accounts of the United States (Z.1) of the Federal Reserve Board, the house price index from the
Federal Housing Finance Agency, the consumer and producer price indexes from the Bureau of Labor
Statistics, the bill and bond yields from the Federal Reserve, and the spot prices of WTI crude oil and
Henry Hub gas from the Energy Information Administration, all as published on FRED \parencite{fredch10}. The
futures prices and inventories are from the EIA \parencite{eiach10}. These are works of US federal agencies,
and \texttt{scripts/pull\_ch10.py} downloads them, cuts them at 30 September 2026 and ships them with the
code; the quarterly series end in the second quarter of 2026 and the inventories in July 2026. US equity
returns are from the data library of \textcite{frenchlib}, in the release built from the CRSP database of
August 2026 and retrieved on 19 October 2026; these files are not redistributed, and
\texttt{scripts/build\_french.py -{}-download} rebuilds them. The building, the loans, the forest, the
concession and the exercise inputs are illustrations.

\printbibliography[heading=subbibliography,title={References for Chapter \thechapter}]
\end{refsection}
